\PassOptionsToPackage{hyphens}{url}
\documentclass[aps,prd,twocolumn,superscriptaddress,nofootinbib]{revtex4-2}

% MiKTeX REVTeX 4.2 needs explicit package loading
\usepackage{amsmath}
\usepackage{amssymb}
\usepackage{amsthm}
\usepackage{booktabs}
\usepackage{tabularx}
\usepackage{xcolor}
\usepackage{graphicx}
\graphicspath{{../figures/}}
\input{glyphtounicode}
\pdfgentounicode=1

\newtheorem{definition}{Definition}
\newtheorem{proposition}{Proposition}
\newtheorem{theorem}{Theorem}
\newtheorem{corollary}{Corollary}
\newtheorem{lemma}{Lemma}
\newtheorem{remark}{Remark}
\newtheorem{axiom}{Axiom}

\usepackage[unicode=true]{hyperref}
\hypersetup{
    pdftitle={The Saturation Theorem: Axiomatic Constraints on Quantum Gravity from Cosmological Relaxation},
    pdfauthor={Lukas Geiger},
    pdfsubject={Axiomatic constraints on one-dimensional cosmological saturation dynamics in quantum gravity},
    pdfkeywords={quantum gravity, saturation dynamics, axiomatic cosmology, Curvature Relaxation Model, renormalization group, tanh saturation},
    colorlinks=true,
    linkcolor=black,
    urlcolor=blue,
    citecolor=black
}
\renewcommand{\acknowledgmentsname}{AI Disclosure}

\begin{document}
\pagenumbering{gobble}
\onecolumngrid

% ===================================================================
% TITLE PAGE
% ===================================================================

\thispagestyle{empty}
\vspace*{0.08\textheight}

\begin{center}
{\LARGE The Saturation Theorem: Axiomatic Constraints on Quantum Gravity\par}
\vspace{0.2em}
{\LARGE from Cosmological Relaxation\par}
\vspace{1.8em}
{\large Lukas Geiger\textsuperscript{1,*}\par}
\vspace{0.6em}
{\itshape \textsuperscript{1}Independent Researcher, Bernau, Germany\textsuperscript{\textdagger}\par}
\vspace{0.2em}
{(Dated: \today)\par}
\end{center}

\vspace{1.8em}
\noindent
We formulate four minimal axioms -- finite geometric capacity, cooperative reinforcement, monotone cosmic control, and renormalization-group composition -- as candidate macroscopic constraints for any quantum gravity theory whose cosmological limit reproduces general relativity with a bounded relaxation mechanism. These axioms imply an Abel-coordinate description of one-dimensional saturation dynamics once the reversible one-dimensional response-channel hypothesis $D^\pm$ is added. A hyperbolic tangent profile follows once the saturation boundary is fixed by the additional projective double-boundary quotient assumption D': the ratio $R(x)=(1+x)/(1-x)$ must multiply under composition. Without that assumption, smooth boundary regularity alone leaves a family of smooth Abel collars, for example Abel coordinates of the form $\operatorname{arctanh}(x)+\varepsilon x$. The corrected \textit{Saturation Theorem} is therefore a projective-collar normal-form theorem: $S(g)=S_{\max}\tanh(\alpha\,u(g))$ is selected by scale coherence plus the projective boundary quotient, not by bare smoothness alone. Major quantum gravity programs -- Loop Quantum Gravity, asymptotic safety, string/holographic approaches, causal set theory, and noncommutative geometry -- provide motivation for A--D when restricted to cosmological observables, while D' remains a separate diagnostic assumption to be physically derived. Under those hypotheses, the $\tanh$ saturation ODE $dX/da = k(1-X^2)$ of the Curvature Relaxation Model can be read as the projective normal form of capacity-limited cooperative relaxation, while leaving open which UV completion selects $k$, $\Phi_0$, and the projective quotient physically.

\vfill
{\small
\noindent\textsuperscript{*} ORCID: \url{https://orcid.org/0009-0005-7296-1534}\par
\noindent\textsuperscript{\textdagger} See the AI Disclosure at the end of the paper for tool-use details. All scientific content is the sole responsibility of the author. Paper~V in the CRM series. Companion papers: \cite{Geiger2026,Geiger2026b,Geiger2026c,Geiger2026d}.\par
}

\clearpage
\twocolumngrid
\pagenumbering{arabic}

\iffalse

% ===================================================================
% TITLE PAGE
% ===================================================================

\title{The Saturation Theorem: Axiomatic Constraints on Quantum Gravity\texorpdfstring{\\}{ }from Cosmological Relaxation}

\author{Lukas Geiger}
\thanks{ORCID: \url{https://orcid.org/0009-0005-7296-1534}}
\affiliation{Independent Researcher, Bernau, Germany}

\date{\today}

\begin{abstract}
We formulate four minimal axioms -- finite geometric capacity, cooperative reinforcement, monotone cosmic control, and renormalization-group composition -- as candidate macroscopic constraints for any quantum gravity theory whose cosmological limit reproduces general relativity with a bounded relaxation mechanism. These axioms imply an Abel-coordinate description of one-dimensional saturation dynamics once the reversible one-dimensional response-channel hypothesis $D^\pm$ is added. A hyperbolic tangent profile follows once the saturation boundary is fixed by the additional projective double-boundary quotient assumption D': the ratio $R(x)=(1+x)/(1-x)$ must multiply under composition. Without that assumption, smooth boundary regularity alone leaves a family of smooth Abel collars, for example Abel coordinates of the form $\operatorname{arctanh}(x)+\varepsilon x$. The corrected \textit{Saturation Theorem} is therefore a projective-collar normal-form theorem: $S(g)=S_{\max}\tanh(\alpha\,u(g))$ is selected by scale coherence plus the projective boundary quotient, not by bare smoothness alone. Major quantum gravity programs -- Loop Quantum Gravity, asymptotic safety, string/holographic approaches, causal set theory, and noncommutative geometry -- provide motivation for A--D when restricted to cosmological observables, while D' remains a separate diagnostic assumption to be physically derived. Under those hypotheses, the $\tanh$ saturation ODE $dX/da = k(1-X^2)$ of the Curvature Relaxation Model can be read as the projective normal form of capacity-limited cooperative relaxation, while leaving open which UV completion selects $k$, $\Phi_0$, and the projective quotient physically.
\end{abstract}

\keywords{quantum gravity, saturation dynamics, axiomatic cosmology, conditional normal-form theorem, Curvature Relaxation Model, tanh normal form, renormalization group}

\thanks{See the AI Disclosure at the end of the paper for tool-use details. All scientific content is the sole responsibility of the author.}

\thanks{Paper~V in the CRM series. Companion papers: \cite{Geiger2026,Geiger2026b,Geiger2026c,Geiger2026d}.}

\begin{titlepage}
\maketitle
\end{titlepage}

\fi


% ===================================================================
% 1. INTRODUCTION
% ===================================================================
\section{Introduction}
\label{sec:intro}

\subsection{The problem}

The unification of quantum mechanics and general relativity remains the central open problem of theoretical physics. After decades of effort, several major programs have produced deep structural insights without converging on a single theory:

\begin{itemize}
\item \textbf{Effective field theory (EFT) of gravity} computes reliable quantum corrections at low energies but provides no UV completion \cite{Donoghue1994}.
\item \textbf{Asymptotic safety} seeks a UV fixed point of the gravitational renormalization group (RG) flow, constraining the theory without new degrees of freedom \cite{Weinberg1979,Reuter1998,Percacci2017}.
\item \textbf{String theory} provides a perturbatively UV-finite framework that includes gravity, but faces the landscape/selection problem \cite{Polchinski1998}.
\item \textbf{Loop Quantum Gravity (LQG)} achieves background independence with discrete geometry operators, but struggles with dynamics and the classical limit \cite{Rovelli2004,Thiemann2007}.
\item \textbf{Holography/AdS-CFT} offers a nonperturbative definition of quantum gravity in anti-de~Sitter space, with unclear extension to realistic cosmology \cite{Maldacena1999}.
\item \textbf{Causal set theory} posits discrete causal structure as fundamental, with the continuum as emergent \cite{Sorkin2003,Surya2019}.
\item \textbf{Noncommutative geometry} reformulates geometry as spectral data, thereby linking operators and spacetime \cite{Connes1994}.
\end{itemize}

Orthogonal to these UV-completion programs, the \textit{Swampland conjectures} \cite{Obied2018} constrain which low-energy theories can be consistently embedded in quantum gravity, while the \textit{EFT of Dark Energy} \cite{Gubitosi2013} parametrizes deviations from $\Lambda$CDM at the perturbative level. Neither framework, however, determines the \textit{functional form} of the cosmological saturation profile.

Several UV-completion programs contain motifs that can be compared with bounded cosmological response: finite scales, monotone branches, coarse-graining, or ordered phases. Calling all of these motifs ``saturation'' or ``cooperation'' is initially a source-to-target analogy, not a shared microscopic mechanism. This paper asks which explicit invariants are sufficient for a conditional normal-form statement.

\subsection{The CRM as empirical anchor}

The Curvature Relaxation Model (CRM) \cite{Geiger2026,Geiger2026b,Geiger2026c,Geiger2026d} provides a concrete, empirically tested cosmological framework in which saturation dynamics plays the central role. The model replaces the cosmological constant $\Lambda$ and particulate dark matter with geometric curvature return, achieving $\Delta\chi^2 = -5.5$ vs.\ $\Lambda$CDM on joint SN+CMB+BAO data with 6 parameters and zero Early Dark Energy \cite{Geiger2026b}. The entire framework derives from a single dynamical equation, here written in the additive normal-form coordinate $u$:
\begin{equation}
\frac{dX}{du} = \alpha(1 - X^2)\,,
\label{eq:saturation_ode}
\end{equation}
where $X = \Omega_\Phi/\Phi_0$ is the normalized curvature return potential. The phenomenological CRM uses the scale factor $a$ as parameter, with solution $X(a) = \tanh(k(a - a_{\mathrm{trans}}))$, providing late-time acceleration (``dark energy'') in its saturated phase and gravitational scaffolding (``dark matter'') in its unsaturated phase \cite{Geiger2026b}. \emph{Coordinate caveat (added in this revision):} if a general parameter $a$ is used instead of the additive normal-form coordinate $u$, then $dX/da = \alpha\,u'(a)(1-X^2)$; the constant-coefficient CRM equation in $a$ therefore requires the \emph{additional} affine-channel assumption $u(a) = \lambda(a - a_{\mathrm{trans}})$. The theorem of this paper establishes the projective normal form in $u$; it does not establish that affine identification for any particular UV completion (see Sec.~\ref{sec:discussion}).

Paper~III \cite{Geiger2026c} historically proposed this ODE through a direct trace-coupled effective Lagrangian. The July~2026 truth audit found that route insufficient as a physical derivation and retained it only as a documented failure audit. The active Route~3 reconstruction instead separates a shift-symmetric dust candidate, a vector--scalar MOND sector, and the saturation scalar; it has not derived the saturation ODE or D' from one common Lagrangian. The theorem below is therefore logically independent of that historical derivation claim.

The present paper recasts this observation from \textit{structural analogy} as a conditional normal-form theorem: the four physically motivated axioms, plus B' and $D^\pm$, give an Abel-coordinate saturation law, and the additional projective double-boundary quotient assumption D' selects the $\tanh$ representative.

\subsection{Scope and structure}

The central result is the \textit{Saturation Theorem} (Theorem~\ref{thm:saturation}): Axioms~A--D, together with signed antisymmetry B' and the separate one-dimensional reversible-channel hypothesis $D^\pm$, reduce the effective cosmological response to an Abel coordinate, while the projective double-boundary quotient assumption D' uniquely selects the $\tanh$ normal form up to reparametrization. The paper is organized as follows:

\begin{itemize}
\item Section~\ref{sec:axioms}: Formulation of the four axioms.
\item Section~\ref{sec:theorem}: Statement and proof of the Saturation Theorem.
\item Section~\ref{sec:exclusion}: Non-projective composition laws rather than profile names.
\item Section~\ref{sec:qg}: Compatibility map with quantum-gravity motifs.
\item Section~\ref{sec:ontology}: Ontological implications: spacetime as an ordered phase.
\item Section~\ref{sec:observational}: Observational consequences and open questions.
\item Section~\ref{sec:discussion}: Discussion and outlook.
\end{itemize}


% ===================================================================
% 2. THE FOUR AXIOMS
% ===================================================================
\section{The Four Axioms}
\label{sec:axioms}

We formulate the axioms in terms of an effective response function $S_T$. To keep the roles distinct, $r\geq0$ denotes a physical UV-to-IR control magnitude or scale ratio, $u$ denotes the additive signed normal-form coordinate used after passing to a one-dimensional composition channel, and $x=S/S_{\max}$ denotes the normalized response. The symbol $g$ is retained below as a generic channel variable when no distinction between $r$ and $u$ is needed.

\subsection{Axiom A: Finite Geometric Capacity}

\begin{axiom}[Finite Geometric Capacity]
\label{ax:capacity}
Per Hubble patch, there exists a finite maximum geometric order $C_{\max} < \infty$. Equivalently, the effective response function $S_T(r)$ on the selected positive physical response channel is globally bounded:
\begin{equation}
0 \leq S_T(r) < S_{\max} < \infty \quad \forall\, r \geq 0\,.
\end{equation}
\end{axiom}

\textit{Physical content.} This axiom imposes finite capacity for the selected scalar response channel. It is not a claim that every microscopic observable or every UV amplitude is finite in the underlying theory. Finite area quanta in LQG \cite{Rovelli2004}, the Bekenstein bound \cite{Bekenstein1981}, finite code subspaces \cite{Almheiri2015}, causal-set discreteness \cite{Sorkin2003}, and fixed-point behavior \cite{Reuter1998} are source-domain motifs only. Each requires an explicit map to the bounded response $S_T$ before it counts as structural support for Axiom~A.

Without finite capacity, the framework would permit unbounded geometric configurations associated with the singular behaviors that quantum-gravity programs typically aim to tame.

\textit{Connection to CRM.} The saturation amplitude $\Phi_0$ in the CRM \cite{Geiger2026} is the physical realization of $S_{\max}$. The P\"oschl-Teller potential $V(\phi) = V_0/\cosh^2(\phi/\phi_0)$ enforces $\Omega_\Phi \leq \Phi_0$ dynamically \cite{Geiger2026c}.

\subsection{Axiom B: Cooperative Reinforcement}

\begin{axiom}[Cooperative Reinforcement]
\label{ax:cooperative}
Existing geometric order facilitates further ordering. The effective response function $S_T(g)$ is smooth ($C^\infty$) for $g \geq 0$, with a local expansion at $g = 0$:
\begin{equation}
S_T(g) = c_1 g + c_3 g^3 + \mathcal{O}(g^5)\,, \quad c_1 > 0\,.
\label{eq:ir_expansion}
\end{equation}
The odd-power structure reflects the catalytic nature of the response: the rate of geometric ordering is proportional to both the existing ordered fraction and the remaining capacity.
\end{axiom}

\textit{Physical content.} Formally, $c_1>0$ preserves only the sign of the local response, while the odd-power structure is consistent with the signed extension (Axiom~B'). The word ``cooperative'' is a motivating analogy to positive feedback in an order parameter; it does not identify microscopic constituents, interactions, correlations, or a ferromagnetic Hamiltonian.

Within this effective model, removing the positive local response would remove the intended self-reinforcing saturation interpretation. It would not prove that classical spacetime is impossible. Spontaneous symmetry breaking is therefore used below as a heuristic source domain whose preserved and unpreserved invariants are stated explicitly.

\textit{Connection to CRM.} The saturation coupling $k$ in the CRM ODE~\eqref{eq:saturation_ode} supplies the positive local coefficient $c_1$. The game-theoretic equilibrium of Paper~I \cite{Geiger2026} motivates the language of cooperative ordering; it is not a microscopic derivation of that mechanism.

A companion condition extends the domain of $S_T$ to negative arguments, enabling the group-theoretic classification of Step~2 in the proof:

\begin{axiom}[Antisymmetric Extension (B')]
\label{ax:antisymmetry}
The response function $S_T$ admits a smooth antisymmetric extension to all $g \in \mathbb{R}$:
\begin{equation}
S_T(-g) = -S_T(g) \quad \forall\, g\,.
\label{eq:antisymmetry}
\end{equation}
\end{axiom}

\textit{Physical content.} The antisymmetry has three complementary motivations:

\textit{(i)~Sign covariance of the response.} The parameter $g$ is not a ``distance'' (which would be non-negative by definition) but a \textit{signed deviation} from a symmetric reference state. In the CRM, $g$ parametrizes the curvature perturbation relative to the de~Sitter background; positive $g$ corresponds to over-dense perturbations that enhance geometric ordering, and negative $g$ to under-dense regions that reduce it. Sign covariance ($S_T(-g) = -S_T(g)$) then states that the response function treats over- and under-dense regions symmetrically in magnitude -- a natural physical requirement absent any symmetry-breaking mechanism in the response itself. We emphasize that Axiom~B' is \textit{not} a claim about cosmological time-reversal symmetry (which is broken by the expansion); it is a constraint on the functional form of the response under $g \to -g$.

\textit{(ii)~Precedents in physics.} Sign covariance of response functions is familiar in symmetric response channels: in electrodynamics, the polarization $P(-E) = -P(E)$ for centrosymmetric media, and in magnetism, $M(-H) = -M(H)$ for the idealized paramagnetic/ferromagnetic response. The odd extension in Axiom~B' is a modelling assumption in a locally signed response coordinate; it is not asserted that generic beta functions are odd.

\textit{(iii)~Mathematical role.} Together with the odd-power structure of Axiom~B, B' extends the composition law from the half-line $[0, S_{\max})$ to a full group operation on $(-S_{\max}, S_{\max})$, which is essential for the Abel-equation classification. Without B', only a semigroup structure is available, and the Abel-representation argument of Step~2 does not apply in the same form. We note that this is an explicit assumption, not a derived consequence; its relaxation leads to alternative composition structures (e.g., $C(x,y) = x + y - xy$ on $[0,1)$), which is why it appears as a separate axiom rather than being subsumed into Axiom~B.

\subsection{Axiom C: Monotone Cosmic Control}

\begin{axiom}[Monotone Cosmic Control]
\label{ax:monotone}
The effective response function $S_T(g)$ is monotonically increasing, with saturation at infinity:
\begin{equation}
S_T'(g) \geq 0\,, \quad \lim_{g \to \infty} S_T(g) = S_{\max}\,.
\label{eq:monotone}
\end{equation}
There exists exactly one relevant direction (control parameter) governing the transition from low to high geometric order.
\end{axiom}

\textit{Physical content.} This excludes oscillating, multi-fixed-point, or chaotic geometric dynamics. It guarantees:
\begin{itemize}
\item A unique time direction (thermodynamic arrow).
\item No ``bouncing'' between geometric phases.
\item Stable expansion toward a de~Sitter attractor.
\end{itemize}

In the CRM, the control parameter is the scale factor $a$; in RG language, it corresponds to the single relevant direction flowing from the UV fixed point to the IR \cite{Percacci2017}. Axiom~C is violated by cyclic cosmologies and by models with multiple stable vacua -- both of which lack a unique thermodynamic arrow. We note that Axiom~C is a restriction to the \textit{cosmological sector} with a single dominant transition: it does not address the string landscape (where tunneling between metastable vacua produces non-monotone trajectories) or stochastic inflation (where local fluctuations break monotonicity). These scenarios involve multiple effective control parameters or stochastic dynamics, placing them outside the scope of the single-field, deterministic axiom system. The theorem should be read as follows: A--D, with B' and the interior group form $D^\pm$, give an Abel/collar class for a single deterministic saturation channel; adding D' selects the projective $\tanh$ representative.

\textit{Connection to CRM.} The monotonicity of $\Omega_\Phi(a)$ is a direct consequence of the ODE~\eqref{eq:saturation_ode}: since $1 - X^2 > 0$ for $|X| < 1$, the solution is strictly increasing.

\begin{remark}[Axiom interdependence]
Lemma~\ref{lem:global_monotonicity} and Corollary~\ref{cor:monotonicity} isolate the exact dependency: the monotonicity and endpoint-saturation clauses of Axiom~C follow only after the reversible response law $D^\pm$ is combined with a genuine global one-parameter control-group homomorphism. The response-side group law alone does not supply that global control map. We therefore retain Axiom~C explicitly in the main theorem for physical transparency and for models in which only a local chart, semigroup, stochastic channel, or multi-field control is known.
\end{remark}

\subsection{Axiom D: Renormalization-Group Composition}

\begin{axiom}[Renormalization-Group Composition]
\label{ax:composition}
There exists a scale composition $g \mapsto g_1 \oplus g_2$ (``two RG steps in succession'') such that the response function is compositionally coherent:
\begin{equation}
S_T(g_1 \oplus g_2) = \mathcal{C}\!\left(S_T(g_1),\, S_T(g_2)\right)\,,
\label{eq:composition}
\end{equation}
where $\mathcal{C}$ is a smooth, commutative, associative composition with neutral element $0$ and absorbing boundary $S_{\max}$:
\begin{equation}
\mathcal{C}(S, 0) = S\,, \quad \mathcal{C}(S, S_{\max}) = S_{\max}\,.
\label{eq:composition_boundary}
\end{equation}
Furthermore, $\mathcal{C}$ extends smoothly (i.e., $C^\infty$) to the closed square $[0, S_{\max}]^2$.
\end{axiom}

\begin{axiom}[One-Dimensional Reversible Channel ($D^\pm$)]
\label{ax:reversible_channel}
After normalization by $S_{\max}$ and extension by B', the induced operation on the signed response interval $C:(-1,1)^2\to(-1,1)$ is cancellative and strictly monotone in each argument, and each translation $y\mapsto C(x,y)$ is a smooth diffeomorphism with inverse $y\mapsto C(-x,y)$. All Abel-coordinate conclusions below use this additional one-dimensional deterministic-channel hypothesis; it is not a consequence of the bare positive-face semigroup law alone.
\end{axiom}

\textit{Notation for channel variables.} Let $r\geq0$ denote the physical UV-to-IR control magnitude or scale ratio, $u$ the additive signed normal-form coordinate with $u(g_1\oplus g_2)=u(g_1)+u(g_2)$, and $x=S/S_{\max}$ the normalized response. Statements about constant coefficients are statements in $u$, not automatically in an externally chosen physical parameter. The positive physical branch is recovered by restricting the signed local channel to the physically realized side.

\begin{remark}[Boundary-collar guardrail]
The smooth-extension clause in Axiom~D is not, by itself, used below as a uniqueness claim for $\mathrm{arctanh}$. The Abel coordinate
\[
  \phi_\varepsilon(x)=\operatorname{arctanh}(x)+\varepsilon x
\]
induces a smooth associative composition on the open interval and gives a smooth positive boundary collar via
\[
  q_\varepsilon(x)=e^{-2\phi_\varepsilon(x)}
  =\frac{1-x}{1+x}e^{-2\varepsilon x}.
\]
Thus bare smoothness of the positive saturation face places the response in an Abel-coordinate collar framework; it does not select a unique projective coordinate. Proposition~\ref{prop:smooth_collar_classification} below classifies the smooth strict members of that framework by a boundary-defining function $q_h$. Boundary smoothness is used as an admissibility condition for the response channel. It is not derived from a UV model in this paper. The $\tanh$ normal form becomes a theorem only after the projective double-boundary quotient D' is fixed explicitly.
\end{remark}

\begin{axiom}[Projective Double-Boundary Quotient Assumption (D')]
\label{ax:projective_quotient}
For the normalized response $x=S/S_{\max}\in(-1,1)$, the projective boundary quotient
\[
  R(x):=\frac{1+x}{1-x}
\]
is multiplicative under composition:
\begin{equation}
  R(C(x,y))=R(x)R(y).
  \label{eq:projective_quotient}
\end{equation}
Equivalently, the signed coordinate $h(x):=\frac12\log R(x)=\operatorname{arctanh}(x)$, rather than the ratio itself, is additive under the one-dimensional response composition.
\end{axiom}

\textit{Physical status of D'.} D' is the projective double-boundary quotient assumption. Once the quotient $R=(1+x)/(1-x)$ is fixed prior to the choice of an Abel coordinate, the velocity-addition law follows algebraically. Thus D' is not derived by the A--D hypotheses; it is the additional projective input whose physical derivation is left open. It stipulates that the physically composed quantity is the marked boundary-capacity odds, whose logarithm is analogous to a log-odds or rapidity coordinate. For any candidate UV program, the diagnostic defect is
\[
  \Delta_{D'}(x,y):=\log R(C(x,y))-\log R(x)-\log R(y),
\]
which must vanish for the exact projective $\tanh$ representative. Establishing $\Delta_{D'}=0$ from a microscopic theory remains an open physical task.

\begin{remark}[Marked-projective naturality audit]
Let $b_-=-1$, $e=0$, and $b_+=1$ denote the two response boundaries and the neutral response. Then
\[
  R(x)=
  \frac{(x-b_-)/(b_+-x)}{(e-b_-)/(b_+-e)}
  =\frac{1+x}{1-x}
\]
is a normalized four-point cross-ratio, and $h(x)=\frac12\log R(x)$ is the signed Hilbert coordinate from $e$ on the interval \cite{PapadopoulosTroyanov2008,RainioVuorinen2023}. This quantity is invariant when a projective change of chart acts on both boundaries, the neutral point, and the response. In a chart that keeps only the two boundaries fixed, an orientation-preserving projective automorphism has $R(Tx)=kR(x)$; fixing $T(e)=e$ forces $k=1$, while reversing the boundary orientation sends $R$ to $R^{-1}$. Thus $R$ is mathematically canonical only for an oriented, marked projective response interval, not for an unmarked smooth interval.

That restriction is physical content rather than a consequence of A--D or $D^\pm$. Indeed, the non-projective collar above has the exact multiplicative Abel balance
\[
  Q_\varepsilon(x):=e^{2\phi_\varepsilon(x)}
  =R(x)e^{2\varepsilon x}
\]
under its induced composition $C_\varepsilon$, although the fixed cross-ratio $R$ generally has $\Delta_{D'}\neq0$. Hence the existence of \emph{some} multiplicative balance is generic and does not select D'. A physical selection of $R$ requires an independently defined observable with two finite distinguishable boundary responses, a neutral identity response, and an exact composition law whose signed Hilbert increments add. No such source-derived identification is established here; D' remains a conditional projective gauge fixing and an empirical defect gate, not a derived physical balance law.
\end{remark}

\begin{remark}[External IR-finite transport gate]
An exact physical near-candidate sharpens this boundary without closing it. For the $\nu=1/3$ fractional-quantum-Hall point contact, Fendley, Ludwig, and Saleur compute the nonperturbative conductance $G(T/T_B)$ and its finite endpoint responses $G/G_0\to0,1$, where $G_0=e^2/(3h)$ \cite{FendleyLudwigSaleur1995}. With the source-fixed scale $t=\log(T/T_B)$ and the endpoint normalization $x=2G/G_0-1$, the quoted laws $G/G_0\sim A e^{4t}$ and $1-G/G_0\sim B e^{-4t/3}$ give
\[
 \frac{dx/dt}{1-x^2}\longrightarrow 2
 \quad\hbox{and}\quad
 \frac{2}{3},
\]
respectively. The coefficient-independent mismatch is a source-asymptotic D' failure. Moreover, the exact object is a momentum-dependent $2\times2$ kink/antikink boundary S matrix, whereas $G$ is a thermodynamic-Bethe-ansatz weighted integral of a transmission probability \cite{GhoshalZamolodchikov1994,FendleyLudwigSaleur1995}. Exact star-product composition of full unitary scattering matrices \cite{KostrykinSchrader2001} therefore does not supply a binary scalar law $C(G_1,G_2)$: phase and channel information discarded by $G$ can change the composed transmission. This audit supplies neither the missing physical D' law nor a cosmological or four-dimensional UV transfer.
\end{remark}

\begin{remark}[Massless RG-flow amplitude near-hit]
The exact massless scattering description of the tricritical-Ising-to-Ising flow provides a sharper algebraic near-hit \cite{FendleySaleurZamolodchikov1993}. For the mixed right--left channel,
\[
 \begin{aligned}
 S_{RL}(\theta)&=-\tanh(\theta/2-i\pi/4),\\
                &=\frac{1+i e^\theta}{1-i e^\theta},\\
 \theta&=\log(s/M^2),
 \end{aligned}
\]
with finite limits $S_{RL}\to+1$ and $-1$. Marking $S_{RL}(0)=i$ gives the exact identity $i(1-S_{RL})/(1+S_{RL})=e^\theta$. However, the mark and the resulting law $S(\theta_1)\mathbin{C}_{\rm ind}S(\theta_2)=S(\theta_1+\theta_2)$ are audit-induced. The source-defined factorization is instead the Yang--Baxter product of pairwise scattering operators in a multiparticle state; it does not define a binary RG-response composition or identify $i$ as its physical identity. Moreover, $S_{RL}$ lies on the complex unit circle, not on the declared real Paper-V response interval, and the source's coupling-to-mass map and TBA crossover do not identify a running scalar coupling with this amplitude. Thus the exact marked quotient does not close physical D' and carries no four-dimensional or cosmological transfer.
\end{remark}

\begin{remark}[Integrable defect-fusion control]
Integrable line-defect fusion supplies a source-defined physical composition control rather than an audit-induced law \cite{HeJiangLiu2026,DelfinoMussardoSimonetti1994b}. For the parity-invariant non-topological Ising defect, fusion is the short-distance limit $ma\to0$ of two defects, and the exact multiple-scattering amplitudes close on the same family with
\[
 g_{I,f}=\frac{g_{I,1}+g_{I,2}}{1+g_{I,1}g_{I,2}/4}.
\]
Thus $x=g_I/2$ obeys $x_f=(x_1+x_2)/(1+x_1x_2)$ and exactly
\[
 \frac{1+x_f}{1-x_f}=\frac{1+x_1}{1-x_1}\frac{1+x_2}{1-x_2}.
\]
This is a genuine positive D' semigroup control. It is not yet the physical Paper-V derivation: the source audits the no-bound-state branch $0\leq g_I\leq2$, so the transparent neutral element $g_I=0$ is an endpoint rather than an interior mark of a full signed reversible response interval. Moreover, $x$ is a static defect-coupling label. The reflection probability and defect $g$-function are physical scalar quantities, but the source does not identify either with $x$ or transfer the same quotient law to them. Finally, $ma$ is a separation and $mR$ a finite-size TBA scale; no beta function $g_I(\mu)$ or homomorphism from additive RG time to fusion is supplied. The source explicitly leaves open when fusion can be interpreted as an RG flow. Hence the exact defect-fusion law sharpens, but does not close, physical D' and carries no four-dimensional or cosmological transfer.
\end{remark}

\textit{Physical content of D.} Axiom~D requires the response to be self-consistent under successive scale transformations. This mirrors the semigroup property of renormalization-group maps \cite{WilsonKogut1974}; it is a structural RG transfer only when the response projection obeys the homomorphism~\eqref{eq:rg_response_homomorphism}.

The composition law~\eqref{eq:composition} encodes the requirement that the theory be \textit{scale-coherent} in the following precise sense: the response function $S_T$ is a homomorphism from the additive structure of the control parameter to the composition algebra of responses; i.e., no preferred scale exists at which this homomorphism property breaks down or changes character. The absorbing boundary condition $\mathcal{C}(S, S_{\max}) = S_{\max}$ reflects the fact that fully saturated geometry cannot be further ordered -- it is a fixed point.

The $C^\infty$-extension requirement has a direct physical interpretation: it demands that the composition law remains well-behaved \textit{even near saturation}. In physical terms, this means that the approach to the geometric capacity limit is smooth and does not introduce new divergences or critical exponents. Three physical arguments support this requirement:

\textit{(a)~No new phase transition at saturation.} A composition that develops singular derivatives at the boundary would signal a phase transition \textit{within} the saturation process itself -- i.e., the saturation mechanism would break down precisely where it is needed most. The smoothness condition excludes such pathological behavior. This is distinct from phase transitions at the \textit{onset} of saturation (which are permitted and correspond to the CRM's transition epoch $a_{\mathrm{trans}}$); the boundary regularity concerns the \textit{fully saturated} regime.

\textit{(b)~Distinction from critical phenomena.} While critical exponents and divergent susceptibilities characterize phase transitions at finite temperature, the theorem's saturation boundary $S \to S_{\max}$ is an \textit{asymptotic} response regime ($g \to \infty$), not a critical point. The ferromagnetic source domain suggests comparison with $T\to0$ rather than $T\to T_c$, but it supplies no binary composition law for magnetizations. Accordingly, boundary smoothness remains an explicit response-channel assumption, not a consequence of ferromagnetic thermodynamics.

\textit{(c)~Response regularity is a model postulate.} Generic quantum-field-theory correlators have short-distance singularities, so vacuum language does not derive $C^\infty$ boundary composition. Interpreting full saturation as a maximally ordered state is heuristic; the required smoothness remains an explicit hypothesis of the selected effective response channel.

\textit{Connection to CRM.} For the phenomenologically chosen $\tanh$ response, the addition identity $\tanh(u_1 + u_2) = (\tanh u_1 + \tanh u_2)/(1 + \tanh u_1 \tanh u_2)$ algebraically defines the required law $\mathcal{C}(x,y) = (x + y)/(1 + xy/S_{\max}^2)$. This internal normal-form identity is not a thermodynamic or microscopic derivation of the physical composition variable.


% ===================================================================
% 3. THE SATURATION THEOREM
% ===================================================================
\section{The Saturation Theorem}
\label{sec:theorem}

\subsection{Statement}

\begin{theorem}[Saturation Theorem, projective-collar form]
\label{thm:saturation}
Let $T$ be a theory whose effective response function $S_T(g)$ satisfies Axioms~A--D, the antisymmetry condition~B', and the separate one-dimensional reversible-channel hypothesis $D^\pm$ of Axiom~\ref{ax:reversible_channel}. Then the normalized response $\sigma(g)=S_T(g)/S_{\max}$ admits an odd Abel coordinate $\phi:(-1,1)\to\mathbb{R}$ such that
\[
  \phi(C(x,y))=\phi(x)+\phi(y).
\]
If, in addition, the projective double-boundary quotient assumption D' holds, then there exist constants $\alpha > 0$, $S_{\max} > 0$, and a strictly monotone additive normal-form coordinate $u = u(g)$ with $u(0) = 0$, such that
\begin{equation}
S_T(g) = S_{\max}\,\tanh\!\left(\alpha\,u(g)\right)\,.
\label{eq:saturation_form}
\end{equation}
The composition law is then determined as the relativistic velocity addition $\mathcal{C}(x,y) = (x+y)/(1+xy)$, up to rescaling by $S_{\max}$. Without D', these hypotheses determine the Abel/collar structure but do not exclude all smooth non-projective collars.
\end{theorem}

Table~\ref{tab:proof_status_ledger} separates the logical layers of the corrected result. This prevents the compatibility discussion in Section~\ref{sec:qg} and the empirical remarks in Section~\ref{sec:observational} from being read as proof of the projective quotient.

\begin{table*}[t]
\caption{Proof-status ledger for the corrected Saturation Theorem.}
\label{tab:proof_status_ledger}
\small
\begin{ruledtabular}
\begin{tabular}{@{}llll@{}}
\parbox[t]{0.14\textwidth}{\textbf{Layer}} & \parbox[t]{0.24\textwidth}{\textbf{Input}} & \parbox[t]{0.22\textwidth}{\textbf{Output}} & \parbox[t]{0.27\textwidth}{\textbf{Status / non-claim}} \\
\parbox[t]{0.14\textwidth}{A--D, B', $D^\pm$} & \parbox[t]{0.24\textwidth}{Bounded, cooperative, signed one-dimensional interior composition} & \parbox[t]{0.22\textwidth}{Odd Abel coordinate and the smooth strict $q_h$ collar family} & \parbox[t]{0.27\textwidth}{Classified within the declared single-channel subclass; not yet a unique $\tanh$ profile.} \\[0.6em]
\parbox[t]{0.14\textwidth}{D'} & \parbox[t]{0.24\textwidth}{Projective double-boundary quotient $R(C)=R(x)R(y)$} & \parbox[t]{0.22\textwidth}{Canonical $\tanh$ representative and velocity-addition law} & \parbox[t]{0.27\textwidth}{Additional projective assumption; treated as diagnostic, not derived here from a UV program.} \\[0.6em]
\parbox[t]{0.14\textwidth}{QG program map} & \parbox[t]{0.24\textwidth}{LQG, asymptotic safety, holography, causal sets, CDT and NCG} & \parbox[t]{0.22\textwidth}{Compatibility pattern for A--D, $D^\pm$ and D'} & \parbox[t]{0.27\textwidth}{Motivation and filter; not a verification that any program satisfies all inputs.} \\[0.6em]
\parbox[t]{0.14\textwidth}{CRM data anchor} & \parbox[t]{0.24\textwidth}{Existing CRM fit and monotone transition evidence} & \parbox[t]{0.22\textwidth}{Empirical motivation for using saturation dynamics} & \parbox[t]{0.27\textwidth}{Not an independent test of this theorem, because the CRM already used a $\tanh$ profile.} \\[0.6em]
\parbox[t]{0.14\textwidth}{Perturbed extensions} & \parbox[t]{0.24\textwidth}{Multi-field, stochastic, EFT or UV--IR matching defects} & \parbox[t]{0.22\textwidth}{Future residual/gap test $r_X/g_X \ll 1$} & \parbox[t]{0.27\textwidth}{Recommended stability extension; not part of the present proof.} \\
\end{tabular}
\end{ruledtabular}
\end{table*}

\subsection{Key Lemmas}

The uniqueness of $\mathrm{arctanh}$ rests on a projective quotient, not on smoothness alone.

\begin{lemma}[Projective quotient selects $\mathrm{arctanh}$]\label{lem:projective_arctanh}
Let $C$ be a continuous, strictly monotone, commutative and associative group law on $(-1,1)$ with neutral element $0$. Suppose $C$ has an Abel coordinate $\phi$,
\[
  \phi(C(x,y))=\phi(x)+\phi(y),
\]
and satisfies the projective quotient law~\eqref{eq:projective_quotient}. Then there exists $c>0$ such that
\[
  \phi(x)=c\,\operatorname{arctanh}(x),
  \qquad
  C(x,y)=\frac{x+y}{1+xy}.
\]
\end{lemma}

\begin{proof}
Set $L(x):=\frac12\log((1+x)/(1-x))=\operatorname{arctanh}(x)$. By D',
\[
  L(C(x,y))=L(x)+L(y).
\]
Thus $F(t):=\phi(\tanh t)$ is a continuous additive function of $t=L(x)$, so $F(t)=ct$ for some $c>0$. Hence $\phi(x)=cL(x)=c\,\operatorname{arctanh}(x)$, and solving $L(C)=L(x)+L(y)$ gives $C(x,y)=(x+y)/(1+xy)$.
\end{proof}

\begin{remark}[Why D' is necessary]
The family $\phi_\varepsilon(x)=\operatorname{arctanh}(x)+\varepsilon x$ is odd, smooth and strictly increasing for small $\varepsilon$, diverges at the boundary, and induces a smooth positive saturation collar through $q_\varepsilon=e^{-2\phi_\varepsilon}$. Hence the old claim that $C^\infty$ boundary regularity alone excludes all non-$\mathrm{arctanh}$ Abel coordinates is too strong. The corrected theorem uses D' as the mathematical and physical normalization that removes this boundary-gauge freedom.
\end{remark}

\begin{proposition}[Classification of smooth strict boundary collars]
\label{prop:smooth_collar_classification}
Let $C\in C^\infty([0,1]^2,[0,1])$ be commutative and associative, with neutral element $0$, absorbing element $1$, and strict increase on the open positive face. Assume its interior is the positive restriction of the cancellative one-dimensional channel $D^\pm$. Define
\[
 a(x):=\partial_2C(x,0),
 \qquad
 \phi(x):=\int_0^x\frac{dt}{a(t)}.
\]
Then $a(1)=0$ is a simple zero: with $\kappa:=-a'(1)$ one has $\kappa>0$. Moreover,
\[
 q(x):=e^{-\kappa\phi(x)}
\]
is the unique decreasing $C^\infty$ boundary-defining diffeomorphism $[0,1]\to[1,0]$ satisfying
\[
 \begin{aligned}
 q(0)&=1, & q(1)&=0,\\
 q'(1)&\neq0, & q(C(x,y))&=q(x)q(y).
 \end{aligned}
\]
Conversely, every such $q$ defines one and only one smooth strict positive collar by
\[
 C_q(x,y)=q^{-1}\!\bigl(q(x)q(y)\bigr).
\]
Relative to $q_0(x)=(1-x)/(1+x)$, the full class is parameterized uniquely as
\[
 \begin{aligned}
 q_h(x)&=q_0(x)e^{-2h(x)}, & h(0)&=0,\\
 1+(1-x^2)h'(x)&>0.&&
 \end{aligned}
\]
with $h\in C^\infty([0,1])$. Compatibility with B' is equivalent to $h$ admitting a smooth odd extension through $0$. Its normalized generator and the conjugating response chart are
\[
 \begin{aligned}
 a_h(x)&=
 \frac{(1+h'(0))(1-x^2)}{1+(1-x^2)h'(x)},\\
 H_h&=q_0^{-1}\circ q_h.
 \end{aligned}
\]
\[
 \begin{aligned}
 H_h(C_h(x,y))
 &=C_0(H_h(x),H_h(y)),\\
 C_0(x,y)&=\frac{x+y}{1+xy}.
 \end{aligned}
\]
Finally, in the fixed marked response chart, D' holds if and only if $h=0$. More precisely,
\[
 \Delta_{D'}(x,y)
 =2\bigl(h(x)+h(y)-h(C_h(x,y))\bigr).
\]
The deterministic ledger is
\path{results/paper5/SMOOTH_COLLAR_CLASSIFICATION_GATE_2026-08-26.json}.
\end{proposition}

\begin{proof}
The positive restriction is a smooth strict Archimedean t-conorm, whose standard continuous representation uses an additive generator unique up to positive scale \cite{MesiarovaZemankova2015}; the signed strict extension is an Abelian group in the corresponding bipolar construction \cite{Grabisch2009}. Here the boundary regularity gives the sharper statement. The identity and $D^\pm$ imply $a(0)=1$ and $a>0$ on $[0,1)$, while absorption gives $a(1)=0$. In Abel time $t=\phi(y)$, the translation $x\mapsto C(x,y)$ is the time-$t$ flow of $a(x)\partial_x$. Its derivative at the fixed boundary is
\[
 \partial_1C(1,y)=e^{a'(1)\phi(y)}.
\]
If $a'(1)=0$, this derivative equals $1$ for every $y<1$. Closed-square smoothness instead makes it converge to $\partial_1C(1,1)=0$, since $C(x,1)=1$. Positivity gives $a'(1)\leq0$, hence $a'(1)<0$.

Writing $a(x)=\kappa(1-x)A(x)$ with $A(1)=1$, the equation $a q'=-\kappa q$ shows that $q/(1-x)$ extends smoothly and positively to $x=1$. Thus $q$ is a boundary diffeomorphism, and Abel additivity gives its multiplicative law. Any second multiplicative boundary-defining coordinate is $q^c$ by the uniqueness of continuous additive generators; a simple zero forces $c=1$. Pulling multiplication back by an arbitrary such $q$ proves the converse.

Since both $q$ and $q_0$ have simple zeros, their quotient is smooth and positive through $x=1$, which gives the displayed unique $h$ parameterization and its monotonicity condition. The odd-extension and conjugacy statements follow directly. If D' holds, $q_0$ is a second multiplicative boundary-defining coordinate and uniqueness gives $q_h=q_0$, hence $h=0$; the converse is immediate. Taking logarithms of $q_h(C_h)=q_h(x)q_h(y)$ gives the displayed defect identity.
\end{proof}

\subsection{Proof of the Saturation Theorem}

The proof proceeds in three steps.

\textit{Step 1: Derivation of the functional equation.}
Define the normalized response $\sigma(g) = S_T(g)/S_{\max} \in [0, 1)$. Axioms~A and~C guarantee that $\sigma$ is bounded and monotonically increasing. Axiom~D [Eq.~\eqref{eq:composition}] requires the existence of a composition:
\begin{equation}
\sigma(g_1 \oplus g_2) = C(\sigma(g_1), \sigma(g_2))\,,
\label{eq:normalized_composition}
\end{equation}
where $C: [0,1) \times [0,1) \to [0,1)$ is commutative, associative, has neutral element $0$, and absorbing boundary $1$.

\textit{Step 2: Reduction to Abel's equation.}
By Axiom~B, $\sigma$ has the local expansion~\eqref{eq:ir_expansion} with $c_1 > 0$. Axiom~B' [Eq.~\eqref{eq:antisymmetry}] provides the odd extension $\sigma(-g)=-\sigma(g)$. Together with the reversible-channel hypothesis $D^\pm$, this gives a smooth cancellative group operation on $(-1,1)$ with absorbing boundaries at $\pm 1$ treated only as limiting faces.

Under $D^\pm$, the associativity condition $C(C(x,y),z) = C(x, C(y,z))$ and commutativity define a \textit{one-parameter continuous group law} on $(-1,1)$ (the Taylor expansion at the origin defines a one-dimensional commutative formal group law in the sense of \cite{Hazewinkel1978}; here the connection is that any smooth, associative, commutative binary operation on a neighborhood of $0 \in \mathbb{R}$ with neutral element $0$ is, by Taylor expansion, a formal group law over $\mathbb{R}$, and conversely, Acz\'el's analytic theorem applies to the convergent realization of this formal structure on the interval $(-1,1)$). We apply Acz\'el's representation theorem \cite{Aczel1966} to the \textit{open interval} $(-1,1)$, where $C$ is smooth, strictly monotone in each argument, and the neutral element $0$ lies in the interior. The absorbing boundary conditions $C(x, \pm 1) = \pm 1$ are \textit{not} part of the domain on which Acz\'el's theorem is applied; they are treated as limiting behavior. Specifically, Acz\'el's theorem yields a continuous, strictly monotone function $\phi: (-1,1) \to \mathbb{R}$ with $\phi(0) = 0$. The absorbing boundary conditions then imply $\phi(x) \to +\infty$ as $x \to 1^-$ and $\phi(x) \to -\infty$ as $x \to -1^+$ (since $C(x,y) \to 1$ as $y \to 1$ forces $\phi(x) + \phi(y) \to \infty$, which requires $\phi(y) \to \infty$):
\begin{equation}
\phi(C(x,y)) = \phi(x) + \phi(y)\,.
\label{eq:abel}
\end{equation}
This is Abel's functional equation. \emph{(Proof step corrected in this revision; an earlier version wrote $\phi'(0)=c_1$ from Axiom~B, which was a category error: $c_1$ is the derivative of the unnormalized response $S_T$ at the control-parameter origin, not of the Abel coordinate $\phi$ at the response origin.)} The clean construction proceeds via the invariant vector field: let $a(x) = \partial_2 C(x,0)$; under $D^\pm$, translations are orientation-preserving diffeomorphisms, hence $a(x) > 0$ on $(-1,1)$, and one may define $\phi(x) = \int_0^x dt/a(t)$. The invariance of $a$ under the group law then gives $\phi(C(x,y)) = \phi(x) + \phi(y)$, and $\phi$ inherits the $C^\infty$ regularity of $C$ on $(-1,1)$ directly from this integral representation. The function $\phi$ is determined up to a positive multiplicative constant, and antisymmetry (Axiom~B') lets it be chosen odd. Axiom~B supplies $d\sigma/dg(0)=c_1/S_{\max}>0$; if an orientation-preserving additive control coordinate $u=u(g)$ is used, then $\kappa:=(d\sigma/du)(0)>0$ enters only through $(\phi\circ\sigma)'(0)=\phi'(0)\kappa$, fixing the overall positive scale $\alpha$. At this stage Axioms~A--D give an Abel/collar theorem, not yet $\phi=\mathrm{arctanh}$.

\textit{Step 2b: Projective double-boundary quotient.}
If the additional assumption D' holds, Lemma~\ref{lem:projective_arctanh} gives $\phi=c\,\mathrm{arctanh}$. D' is the projective double-boundary quotient assumption, not a consequence of A--D: once $R=(1+x)/(1-x)$ is fixed before choosing an Abel coordinate, the velocity-addition law follows algebraically. This yields the composition law:
\begin{equation}
C(x, y) = \frac{x + y}{1 + x\,y}\,,
\label{eq:addition_formula}
\end{equation}
which is the relativistic velocity addition formula. The uniqueness is therefore projective: it is uniqueness after fixing the boundary quotient $R=(1+x)/(1-x)$, not uniqueness from bare smoothness of all Abel collars.

\textit{Step 3: Inversion to $\tanh$.}
Since $\phi(\sigma(g)) = c\,\mathrm{arctanh}(\sigma(g))$ must be an additive function of the control parameter (i.e., $\phi(\sigma(g)) = \alpha\,u(g)$ for a reparametrization $u$ and constant $\alpha > 0$), we obtain:
\begin{equation}
\sigma(g) = \tanh(\alpha\,u(g))\,,
\end{equation}
and therefore $S_T(g) = S_{\max}\,\tanh(\alpha\,u(g))$, after absorbing the positive constant $c$ into $\alpha u$. When Axioms~A--D, B', $D^\pm$, and D' hold exactly, the result is exact. Approximate axiom satisfaction (e.g., breakdown of the composition law at very small $g$ where standard GR dominates) would introduce corrections; a formal stability analysis of such perturbations is an open problem. \hfill $\square$

\begin{lemma}[Global group monotonicity without a profile premise]
\label{lem:global_monotonicity}
Let $C:(-1,1)^2\to(-1,1)$ be a $C^1$ group law with neutral element $0$, and suppose every translation $L_x(y):=C(x,y)$ is a $C^1$ diffeomorphism. Let
\[
  \sigma:(\mathbb{R},+)\longrightarrow((-1,1),C)
\]
be a $C^1$ global homomorphism with $\sigma(0)=0$ and $\kappa:=\sigma'(0)>0$. Define $a(x):=\partial_2C(x,0)$. Then
\begin{equation}
  a(x)>0,
  \qquad
  \sigma'(u)=\kappa\,a(\sigma(u))>0
  \quad\text{for every }u\in\mathbb{R}.
  \label{eq:global_monotonicity_generator}
\end{equation}
Moreover, $\sigma$ is a strictly increasing bijection from $\mathbb{R}$ onto $(-1,1)$ and therefore $\lim_{u\to\pm\infty}\sigma(u)=\pm1$.
\end{lemma}

\begin{proof}
Because $L_x$ and its inverse are $C^1$, $a(x)=L_x'(0)$ never vanishes. The function $a$ is continuous, $(-1,1)$ is connected, and $a(0)=1$ follows from $C(0,y)=y$; hence $a(x)>0$ throughout the interval. Differentiate the global homomorphism identity
\[
  \sigma(u+v)=C(\sigma(u),\sigma(v))
\]
with respect to $v$ at $v=0$. This gives Eq.~\eqref{eq:global_monotonicity_generator} directly, so strict increase is a conclusion rather than a premise. Since $\sigma'(u)>0$, its image is an open subgroup of the connected target group $((-1,1),C)$; a connected group has no proper open subgroup, so the image is all of $(-1,1)$. A strictly increasing surjection from $\mathbb{R}$ to $(-1,1)$ has the stated endpoint limits. No monotonicity of $\sigma$, Abel coordinate, or projective assumption D' is used in the argument.
\end{proof}

\begin{corollary}[Exact monotonicity dependency for Axiom C]
\label{cor:monotonicity}
For the normalized one-dimensional response, Axiom~B gives $d\sigma/dg(0)=c_1/S_{\max}>0$. If the control admits an orientation-preserving global additive coordinate $u=u(g)$ and the response is a homomorphism $\sigma:(\mathbb{R},+)\to((-1,1),C)$ for all $u,v\in\mathbb{R}$, then $\kappa=(d\sigma/du)(0)>0$ and Lemma~\ref{lem:global_monotonicity} derives the monotonicity and endpoint-saturation clauses of Axiom~C. The single relevant direction is encoded in the one-dimensional global control group; it is not an additional output of the derivative calculation.
\end{corollary}

\begin{remark}
This is a conditional redundancy result for the strengthened global single-channel subclass. The response-side hypothesis $D^\pm$, a local additive chart, or a positive local slope by itself does not establish the global homomorphism gate. The main theorem therefore keeps Axiom~C explicit whenever that gate has not been proved for the physical control variable.
\end{remark}

\begin{remark}[Selberg as a hyperbolic Lie-origin positive control]
For a compact hyperbolic surface $M=\Gamma\backslash\mathbb{H}^2\simeq\Gamma\backslash\mathrm{PSL}(2,\mathbb{R})/\mathrm{SO}(2)$, the Selberg trace-formula setting has a classical source-side operator structure: the Laplace--Casimir operator commutes with the spherical, bi-$K$-invariant convolution channel \cite{Selberg1956}. The projective CRM law is instead the collinear one-dimensional Einstein--Möbius addition. Writing $x=\tanh u$ realizes the boost subgroup $B(u)\in\mathrm{SO}^{+}(1,1)$, for which $B(u)B(v)=B(u+v)$ and therefore $C(x,y)=(x+y)/(1+xy)$ \cite{Ungar2007}. This is a controlled shared rank-one hyperbolic Lie motif; it is neither an identification of the Selberg surface with the response interval nor a derivation of D'.

The stronger ``E10'' transfer proposed in an earlier Selberg idea annex does not follow. Finite hyperbolic area is not finite scalar-response capacity; Casimir centrality is not cooperative reinforcement; antisymmetry of a Lie bracket is not oddness of a response; and geodesic-flow or convolution composition does not supply the response homomorphism required by Axiom~D. No monotone cosmological control, marked saturation boundaries, or projective response observable is defined by the Selberg construction. Selberg is therefore a positive control only for the availability of a Lie-origin invariant operator channel. It supplies no cosmology statement, no UV completion, no full A--D/B' verification, and no strengthening of the Saturation Theorem.
\end{remark}

\begin{proposition}[Axiom non-redundancy]
\label{prop:independence}
The following examples are \emph{separation sketches}, not a complete independence theorem; full pairwise independence of the axiom set is not claimed. \emph{(Two examples corrected in this revision after external review.)} Since Axiom~D is formulated in the bounded response domain, item (1) should be read as an algebraic interior-composition example after the capacity bound is removed, not as a full model of the finite-capacity axiom system.

\textit{(1) Violates A, satisfies B/B' and the algebraic interior-composition part of D:} $\sigma(g) = g$ (unbounded linear response; $C(x,y) = x+y$, associative and commutative).

\textit{(2) Violates the local oddness of B/B', satisfies A,C,D:} $\sigma(g) = 1 - e^{-g}$ for $g \geq 0$, with $C(x,y) = x + y - xy$ (probabilistic or). \emph{Correction:} an earlier version listed this as ``violates B$'$, satisfies B''; in fact the Taylor expansion $1-e^{-g} = g - g^2/2 + g^3/6 - \cdots$ contains a quadratic term, so the example violates the odd local expansion required by Axiom~B itself, not only the signed extension B$'$. It separates $\{$B,B$'\}$ from $\{$A,C,D$\}$ but does not separate B from B$'$.

\textit{(3) Violates the monotonicity clause of C, satisfies A, B/B', and endpoint saturation, but not $D^\pm$:} define the normalized standard bump
\[
 b(z)=\begin{cases}\exp\!\left(1-\dfrac{1}{1-z^2}\right),& |z|<1,\\[0.3em]0,&|z|\geq1.\end{cases}
\]
Then set
\[
 \begin{aligned}
 \psi(g)&=b(4(g-2))-b(4(g+2)),\\
 \sigma(g)&=\tanh(g)+\frac{1}{50}\psi(g).
 \end{aligned}
\]
The perturbation is smooth, odd, supported on $[-9/4,-7/4]\cup[7/4,9/4]$, and vanishes near the origin, so $\sigma(g)=g-g^3/3+O(g^5)$ locally. On $g\geq0$, the maximum ratio of the positive bump to the remaining headroom $1-\tanh g$ is $0.588751887<1$; hence $0\leq\sigma(g)<1$, and compact support preserves $\lim_{g\to\infty}\sigma(g)=1$. Nevertheless,
\[
 \begin{aligned}
 \sigma'\!\left(\frac{17}{8}\right)
 &=\operatorname{sech}^2\!\left(\frac{17}{8}\right)
   -\frac{32}{225}e^{-1/3}\\
 &=-0.046443310\ldots<0.
 \end{aligned}
\]
Thus A, B, B', and endpoint saturation alone do not imply C monotonicity. The deterministic certificate is recorded in \path{results/paper5/AXIOM_C_BUMP_SEPARATION_GATE_2026-08-26.json}. This is not a witness against the joint set $\{A,B,B',D^\pm\}$: the nonmonotone response is noninjective and cannot satisfy the global cancellative response homomorphism, while Corollary~\ref{cor:monotonicity} proves that this stronger gate implies monotonicity. C is retained for physical clarity and to make the single-channel arrow explicit.

\textit{(4) Violates D, satisfies A--C,B':} $\sigma(g) = (2/\pi)\arctan(g)$; bounded, monotone, antisymmetric, but its induced composition (via $\phi = \tan(\pi x/2)$) develops divergent second derivatives ($\sim \epsilon^{-1}$) at the boundary of $[-1,1]^2$, failing the $C^\infty$-extension requirement of Axiom~D. Similarly, $\sigma(g) = g/\sqrt{1+g^2}$ (with $\phi(x) = x/\sqrt{1-x^2}$) has divergent second derivatives at the boundary ($\sim \epsilon^{-1/2}$).
\end{proposition}

\subsection{Interpretation}

The theorem states that $\tanh$ is not a mere \textit{model choice}, but the \textit{canonical projective representative} of a one-dimensional Abel saturation channel. The key mathematical insight is two-step: Axiom~D forces the response function to be the inverse of an Abel coordinate, and D' fixes the otherwise free boundary collar by making the projective quotient $R=(1+x)/(1-x)$ multiplicative.

Physically, the theorem says:
\begin{quote}
\textit{Any quantum gravity theory that is scale-coherent, cooperative, bounded, monotone, signed, reversible in a one-dimensional response channel, and whose saturation boundary satisfies the projective double-boundary quotient assumption D' must produce cosmological saturation dynamics of the $\tanh$ type. Without that quotient assumption, the theorem gives an Abel/collar class rather than a unique profile.}
\end{quote}

The reparametrization $u(g)$ absorbs the microscopic details of the theory: different UV completions correspond to different functions $u(g)$. The common $\tanh(\alpha\,u)$ form is forced only after the projective boundary ratio has been identified as the physical composition variable.


% ===================================================================
% 4. EXCLUSION OF ALTERNATIVES
% ===================================================================
\section{Non-Projective Composition Laws}
\label{sec:exclusion}

\begin{corollary}[Non-projective composition laws rather than profile names]
\label{cor:exclusion}
The exclusions below are exclusions of the induced response composition and boundary quotient under the chosen additive coordinate. They are not coordinate-free exclusions of graph shapes under arbitrary monotone reparametrization. None of the following induced laws can simultaneously satisfy the theorem's structural hypotheses plus the projective quotient assumption D', except when reducible to the $\tanh$ form via admissible reparametrization $u(g)$:
\begin{enumerate}
\item[(i)] \textbf{Logistic profile:} $S_{\max}/(1 + e^{-\beta u})$.
\item[(ii)] \textbf{Rational profile:} $S_{\max}\,u/(1 + u)$.
\item[(iii)] \textbf{Arctangent profile:} $S_{\max}\,(2/\pi)\arctan(\beta u)$.
\item[(iv)] \textbf{Hard saturation:} $S_{\max}\,\min(u, 1)$.
\end{enumerate}
\end{corollary}

\textit{Proof.} Each alternative fails a specific axiom:

\textit{(i) Logistic.} The logistic function $L(u) = 1/(1 + e^{-\beta u})$ has $L(0) = 1/2 \neq 0$, violating the neutral element condition. The shifted logistic $\tilde{L}(u) = (L(u) - 1/2) \cdot 2$ yields $\tilde{L}(u) = \tanh(\beta u/2)$ -- precisely the $\tanh$ form. Therefore, any ``logistic'' model that satisfies Axiom~D is actually $\tanh$ in disguise.

\textit{(ii) Rational.} The function $R(u) = u/(1+u)$ has Abel transform $\phi(x) = x/(1-x)$, yielding $C(x,y) = \phi^{-1}(\phi(x) + \phi(y)) = (x + y - 2xy)/(1 - xy)$. For small $x, y$, this gives $C(x,y) \approx x + y - 2xy$. The quadratic cross-term $-2xy$ is incompatible with the antisymmetry condition (Axiom~B'): an antisymmetric $\sigma$ implies that the induced composition must satisfy $C(-x,-y) = -C(x,y)$, but $(x+y-2xy)/(1-xy)$ does not. Equivalently, $R(u)$ itself is not antisymmetric, violating Axiom~B'.

\textit{(iii) Arctangent.} The function $A(u) = (2/\pi)\arctan(\beta u)$ has Abel transform $\phi(x) = \tan(\pi x/2)$, yielding $C(x,y) = (2/\pi)\arctan(\tan(\pi x/2) + \tan(\pi y/2))$. While this composition is well-defined and smooth on the open square $(-1,1)^2$, it cannot be extended smoothly to the closed square $[-1,1]^2$: the second partial derivatives $\partial^2 C/\partial x^2$ and $\partial^2 C/\partial x\,\partial y$ both diverge as $(x,y) \to (1,1)$ (numerically $\sim \epsilon^{-1}$ along $x = y = 1 - \epsilon$), violating the $C^\infty$ boundary regularity requirement of Axiom~D. This algebraic ($\epsilon^{-1}$) divergence is characteristic of Abel functions with power-law boundary behavior. Logarithmic collars are the admissible boundary class at this face; D' then selects the projective logarithmic collar $\mathrm{arctanh}$ from that class.

\textit{(iv) Hard saturation.} $\min(u, 1)$ is not $C^\infty$ at $u = 1$, violating the smoothness requirement of Axiom~B. \hfill $\square$

\begin{remark}
The logistic case (i) is instructive: it appears to be a genuine alternative, but the composition law plus the projective quotient force it to reduce to $\tanh$. This illustrates the power of scale coherence together with D': surface-level freedom in the profile choice is reduced to a boundary-collar choice, and the projective quotient fixes the CRM rapidity representative.
\end{remark}


% ===================================================================
% 5. COMPATIBILITY MAP WITH QG PROGRAMS
% ===================================================================
\section{Compatibility Map with Quantum-Gravity Motifs}
\label{sec:qg}

We now map how A--D are naturally compatible with motifs in major quantum gravity programs when restricted to cosmological observables. The stronger interior form $D^\pm$ and the projective double-boundary quotient assumption D' are tracked separately as open diagnostics.

\subsection{Loop Quantum Gravity / Loop Quantum Cosmology}

In Loop Quantum Cosmology (LQC) \cite{Ashtekar2006,Bojowald2008}, the modified Friedmann equation reads:
\begin{equation}
H^2 = \frac{8\pi G}{3}\,\rho\!\left(1 - \frac{\rho}{\rho_c}\right)\,,
\label{eq:lqc_friedmann}
\end{equation}
where $\rho_c \sim \rho_{\mathrm{Pl}}$ is a critical density set by the Barbero-Immirzi parameter $\gamma_{\mathrm{BI}}$. The equation has a bounded quadratic response, which is a saturation motif; it is not the CRM response ODE or a shared microscopic capacity law.

\textit{Axiom A:} The critical density $\rho_c$ is a finite-bound motif, but identifying it with the capacity of the selected CRM response requires a map. \textit{Axiom B:} Recovery of the standard Friedmann equation at $\rho\ll\rho_c$ does not by itself establish cooperative reinforcement or B'. \textit{Axiom C:} The expanding branch supplies a monotone segment. \textit{Axiom D:} LQC does not possess a manifest RG composition law in the sense of Axiom~D; a rigorous response homomorphism from the LQG dynamics to Eq.~\eqref{eq:composition} remains open.

Defining $X = 1 - 2\rho/\rho_c$, the LQC Friedmann equation~\eqref{eq:lqc_friedmann} can be rewritten as $H^2 \propto (1 - X^2)$. This preserves a polynomial factor and a finite bound, but the LQC equation governs $H^2$ whereas the CRM ODE governs $dX/da$; no homomorphism between their dynamics is supplied. A proposed parameter correspondence involving the LQG area gap $\Delta$, $\gamma_{\mathrm{BI}}$, $k$, and $\Phi_0$ therefore remains heuristic.

\subsection{Asymptotic Safety}

In the asymptotic safety program \cite{Weinberg1979,Reuter1998,Percacci2017}, the gravitational effective action has a UV fixed point $g^* > 0$ where the dimensionless Newton's constant and cosmological constant become scale-independent. The RG flow from UV to IR defines the effective gravitational coupling $G_k$ and cosmological ``constant'' $\Lambda_k$ as functions of the RG scale $k$.

\textit{Axiom A:} The UV fixed point is a finite-scale motif, not automatically a global bound on the chosen cosmological response. \textit{Axiom B:} A positive critical exponent defines a relevant direction; it does not establish microscopic cooperation. \textit{Axiom C:} Relevance alone does not prove global monotonicity of a projected observable. \textit{Axiom D:} The Wilson RG semigroup $R_{k_1}\circ R_{k_2}=R_{k_1k_2}$ operates on theory space \cite{WilsonKogut1974}, not directly on $S_T$. A structural transfer requires an explicit map $\Pi$ satisfying
\begin{equation}
\Pi(R_{k_1}\circ R_{k_2})=\mathcal{C}\!\left(\Pi(R_{k_1}),\Pi(R_{k_2})\right).
\label{eq:rg_response_homomorphism}
\end{equation}
Restricting to a dominant relevant direction is not sufficient by itself: an invariant one-dimensional channel and the homomorphism~\eqref{eq:rg_response_homomorphism} must still be demonstrated. Neither is derived here for an external UV program.

Asymptotic safety therefore supplies one of the closest source-domain motifs for Axiom~D, while the target-response homomorphism remains open.\footnote{The self-similar, scale-recursive character of Axiom~D is not
confined to cosmological RG flow; the same motif surfaces in the distribution
of the prime numbers. In the primon (Riemann) gas of Julia~\cite{Julia1990}, a
system of non-interacting bosonic oscillators with energies $\log p$ over the
primes $p$ has the Riemann zeta function as its partition function.
Hartnoll and Yang~\cite{HartnollYang2025} have shown that the
Belinsky-Khalatnikov-Lifshitz dynamics of gravity near a spacelike singularity
is, upon semiclassical quantisation, a modular-invariant conformal quantum
mechanics whose states realise exactly such a prime-labelled oscillator gas,
with wavefunctions proportional to automorphic $L$-functions on the critical
axis. This suggests that black-hole singularities may be a further physical
setting in which a capacity-limited, self-similar geometric ordering of the
type encoded by the axioms arises naturally. We record the connection as a
motivating parallel, not as an independent verification of Axioms~A--D.} The asymptotic safety scenario for cosmology \cite{Bonanno2002} produces bounded high-curvature behavior. That is a saturation motif, not evidence for the Paper-V $\tanh$ response or D'.

Recent computations within increasingly sophisticated truncations provide numerical evidence for the Reuter fixed point. Eichhorn and Held~\cite{EichhornHeld2018} report a constraint on the top-quark Yukawa coupling, and Eichhorn~\cite{Eichhorn2020} reviews fixed-point constraints on Standard-Model matter. These results support the RG source domain, but they do not verify the scalar-response capacity, cooperative feedback, global monotonicity, or the homomorphism~\eqref{eq:rg_response_homomorphism} required by Axioms~A--D.

\subsection{String Theory and Holography}

In string theory, UV-finite amplitudes motivate a finite-behaviour analogy, but they do not provide a finite capacity for the selected scalar response. A low-energy EFT within a compactification has an RG flow; this alone does not establish Axioms~B--D or the response homomorphism~\eqref{eq:rg_response_homomorphism}.

The holographic approach (AdS/CFT \cite{Maldacena1999}) relates descriptions across scales, while the Ryu--Takayanagi formula \cite{Ryu2006} relates entanglement entropy and geometric area. These are source-domain motifs for scale relation and monotonicity, not an explicit binary composition on the Paper-V response.

For cosmological applications, the de~Sitter entropy bound $S_{\mathrm{dS}} = \pi/(\Lambda G) < \infty$ and holographic RG flow \cite{deBoer2000} provide finite-bound and scale-flow motifs. Mapping either to $S_T$ and $\mathcal C$ remains an open derivation.

\subsection{Causal Set Theory}

In causal set theory \cite{Sorkin2003,Surya2019}, the cosmological constant arises as a dynamical quantity $\Lambda \sim 1/\sqrt{N}$, where $N$ is the number of causal set elements. Its monotone decrease is a possible control-coordinate motif; it is not already the increasing bounded response required by Axiom~C.

\textit{Axiom A:} Finite sprinkling density is a discreteness motif, not yet the capacity of a scalar response. \textit{Axiom B:} Poisson sprinkling is uncorrelated, and the required positive-feedback invariant has not been derived from sequential growth \cite{Surya2019}. \textit{Axiom C:} The $1/\sqrt{N}$ scaling is monotone, but mapping its decreasing orientation to $S_T$ requires a chosen reparametrization. \textit{Axiom D:} No homomorphism from sequential growth to the response composition has been established.

\subsection{Causal Dynamical Triangulations}

Causal Dynamical Triangulations (CDT) \cite{Ambjorn2005,Loll2020} construct quantum gravity via a regularized path integral over causal triangulations. Monte Carlo simulations reveal a phase diagram with a physically relevant ``de~Sitter phase'' (phase~$C_{dS}$) in which the emergent geometry has Hausdorff dimension $d_H \approx 4$ and a volume profile well-described by Euclidean de~Sitter space.

\textit{Axiom A:} A finite triangulation is a regulator-level finite-count motif; a continuum scalar capacity has not been derived. \textit{Axiom B:} Growth of the de~Sitter volume profile does not by itself establish cooperative feedback or B'. \textit{Axiom C:} The expanding branch supplies a monotone segment. \textit{Axiom D:} Blocking transformations are a coarse-graining motif, but no associative response homomorphism has been established.

\subsection{Noncommutative Geometry}

Connes' program \cite{Connes1994} reconstructs geometry from spectral data (the Dirac operator on a noncommutative algebra). The spectral action principle yields the Standard Model Lagrangian coupled to gravity from the spectrum of a single operator.

\textit{Axiom A:} The spectral cutoff $\Lambda$ is not by itself a finite capacity of the selected response. \textit{Axiom B:} No signed odd response channel has been derived from the heat-kernel expansion. \textit{Axiom C:} A monotone scalar projection has not been specified. \textit{Axiom D:} Products of spectral triples do not by themselves define the binary response law $\mathcal C$. These remain source-domain motifs rather than structural transfers.

\subsection{Summary: Why Conditional Abel/Collar Universality?}

Table~\ref{tab:axiom_check} summarizes the compatibility map. No single external QG program has been rigorously proven here to satisfy A--D, $D^\pm$, and D' simultaneously. The entries record source-domain motifs and open mapping obligations, not proofs of transfer. The theorem establishes universality only \emph{inside} the explicitly defined class of one-dimensional responses satisfying its hypotheses: A--D, B', and $D^\pm$ give an Abel/collar class, while D' selects its projective $\tanh$ representative. This conditional normal-form equivalence is not a Wilsonian universality class and does not imply shared microscopic dynamics or critical exponents \cite{WilsonKogut1974}.

\begin{table*}[t]
\caption{Compatibility map with quantum-gravity motifs. $\checkmark$ = compatible motif with explicit construction, ($\checkmark$) = compatible motif requiring mild assumptions, ? = open derivation, no = known obstruction. The D' column records whether the projective double-boundary quotient is derived or explicitly assumed.}
\label{tab:axiom_check}
\begin{ruledtabular}
\begin{tabular}{lccccc}
Program & A & B/B' & C & D & D' \\
\hline
LQG/LQC & ($\checkmark$) & ? & ($\checkmark$) & ? & ? \\
Asymptotic safety & ($\checkmark$) & ? & ? & ($\checkmark$) & ? \\
Strings/holography & ($\checkmark$) & ? & ? & ? & ? \\
Causal sets & ($\checkmark$) & ? & ($\checkmark$) & ? & ? \\
CDT & ($\checkmark$) & ? & ($\checkmark$) & ? & ? \\
NCG & ? & ? & ? & ? & ? \\
CRM (effective) & $\checkmark$ & $\checkmark$ & $\checkmark$ & $\checkmark$ & assumed \\
\end{tabular}
\end{ruledtabular}
\end{table*}


% ===================================================================
% 6. ONTOLOGY: ORDER, NOT METRIC
% ===================================================================
\section{Spacetime as an Ordered Phase}
\label{sec:ontology}

The compatibility motifs in Table~\ref{tab:axiom_check} raise a conceptual question: if an external response channel is eventually shown to satisfy the stated gates, how should that conditional normal-form equivalence be interpreted without inferring a shared microscopic ontology?

\subsection{The conceptual shift}

The Saturation Theorem, combined with the CRM's empirical success, suggests a possible interpretive reading of quantum gravity:

\begin{quote}
\textit{In this reading, quantum gravity is not merely the quantization of the metric, but also the theory of how geometric order emerges from a pre-geometric quantum system.}
\end{quote}

The $\tanh$ saturation is therefore best read as a conditional signature of a capacity-limited cooperative ordering process, once the one-dimensional composition law and D' are fixed. This does not make it independent evidence for a particular microscopic ontology; it supplies a common normal-form language for spin-network, string/holographic, causal-set, or code-subspace realizations when they satisfy the stated gates.

\subsection{The ferromagnetic analogy}

The analogy with spontaneous magnetization is useful at the level of a mean-field Ginzburg-Landau free-energy model. The CRM saturation ODE $dX/da = k(1 - X^2)$ can be written as the equation of motion for the order parameter of a second-order phase-transition toy model with double-well free energy $F(X) = -k(X - X^3/3)$. Paper~III \cite{Geiger2026c} develops this analogy in detail, identifying:

\begin{center}
\begin{tabular}{ll}
\hline
\shortstack[l]{Ferromagnetic\\source domain} & \shortstack[l]{CRM target\\domain} \\
\hline
Magnetization $M$ & Curvature return $\Omega_\Phi$ \\
\shortstack[l]{Inverse temperature\\$J/k_BT$} & \shortstack[l]{Control coordinate\\$a-a_{\mathrm{trans}}$\\(heuristic reparametrization)} \\
Curie point ($J/k_BT_c$) & Transition $a_{\mathrm{trans}}$ \\
Spin interaction $J$ & Curvature coupling $k$ \\
Saturation $M_s$ & Saturation $\Phi_0$ \\
$\tanh(J/k_BT)$ & $\tanh(k(a - a_{\mathrm{trans}}))$ \\
\hline
\end{tabular}
\end{center}

\begin{table*}[t]
\caption{Invariant audit for the source-to-target transfers used in Paper~V. ``Structural'' is reserved for an explicit normal-form map or homomorphism; a shared word or graph shape is insufficient.}
\label{tab:metaphor_invariants}
\begin{ruledtabular}
\begin{tabular}{llll}
\parbox[t]{0.14\textwidth}{Transfer} & \parbox[t]{0.25\textwidth}{Invariant retained} & \parbox[t]{0.32\textwidth}{Not retained or not derived} & \parbox[t]{0.18\textwidth}{Admissible status} \\
\hline
\parbox[t]{0.14\textwidth}{Saturation} & \parbox[t]{0.25\textwidth}{Normalized bounded interval, ordered endpoints, monotone endpoint approach} & \parbox[t]{0.32\textwidth}{Shared microscopic equation, free energy, or UV mechanism} & \parbox[t]{0.18\textwidth}{Structural only at the normalized response/normal-form level} \\[0.5em]
\parbox[t]{0.14\textwidth}{Ferromagnetic} & \parbox[t]{0.25\textwidth}{Signed scalar order parameter and a mean-field $\tanh$-shaped response} & \parbox[t]{0.32\textwidth}{Magnetization composition law, lattice correlations, critical exponents, constituents} & \parbox[t]{0.18\textwidth}{Heuristic analogy} \\[0.5em]
\parbox[t]{0.14\textwidth}{Cooperative} & \parbox[t]{0.25\textwidth}{Positive local response $c_1>0$} & \parbox[t]{0.32\textwidth}{Spin-like interaction, many-body feedback, correlation function} & \parbox[t]{0.18\textwidth}{Local formal correspondence} \\[0.5em]
\parbox[t]{0.14\textwidth}{Renormalization} & \parbox[t]{0.25\textwidth}{Associativity of successive scale maps if Eq.~\eqref{eq:rg_response_homomorphism} holds} & \parbox[t]{0.32\textwidth}{Projection from theory space to the scalar response for any external QG program} & \parbox[t]{0.18\textwidth}{Conditional structural gate} \\[0.5em]
\parbox[t]{0.14\textwidth}{Temperature} & \parbox[t]{0.25\textwidth}{Order of a chosen one-dimensional control after reparametrization} & \parbox[t]{0.32\textwidth}{Units, thermal ensemble, conjugacy, or physical identification with $a$} & \parbox[t]{0.18\textwidth}{Heuristic reparametrization} \\[0.5em]
\parbox[t]{0.14\textwidth}{Capacity} & \parbox[t]{0.25\textwidth}{Finite response endpoint under $M_s\leftrightarrow\Phi_0$} & \parbox[t]{0.32\textwidth}{Entropy, information, code, or holographic capacity} & \parbox[t]{0.18\textwidth}{Formal bound; physical reading heuristic} \\
\end{tabular}
\end{ruledtabular}
\end{table*}

Mean-field Ising equilibrium obeys a self-consistency equation of the form $m=\tanh[\beta(B+Jqm)]$; it does not define an Abel binary composition of magnetizations. The appearance of $\tanh$ is therefore not by itself a homomorphism. Exact two-dimensional magnetization also has non-mean-field critical behavior \cite{Yang1952}, and RG universality concerns fixed points and scaling data rather than a shared sigmoid \cite{WilsonKogut1974}. Table~\ref{tab:metaphor_invariants} consequently keeps the explicit normalized response/normal-form map as structural, while ferromagnetism, temperature, and microscopic cooperation remain heuristic unless their missing source-to-target maps are derived.

Having established the ontological framework, we now turn to its empirical consequences.

% ===================================================================
% 7. OBSERVATIONAL CONSEQUENCES
% ===================================================================
\section{Observational Consequences and Open Questions}
\label{sec:observational}

\subsection{What the theorem predicts}

The Saturation Theorem has two classes of test targets. Table~\ref{tab:observational_ledger} separates theorem-facing tests from CRM-specific or UV-dependent evidence.

\begin{table*}[t]
\caption{Observational and UV-evidence ledger. The table separates what is tested by the normal-form theorem from what belongs to the CRM fit or to a future UV completion.}
\label{tab:observational_ledger}
\small
\begin{ruledtabular}
\begin{tabular}{@{}llll@{}}
\parbox[t]{0.14\textwidth}{\textbf{Target}} & \parbox[t]{0.24\textwidth}{\textbf{Current role}} & \parbox[t]{0.25\textwidth}{\textbf{Claim status}} & \parbox[t]{0.28\textwidth}{\textbf{Next gate}} \\
\parbox[t]{0.14\textwidth}{Profile fit} & \parbox[t]{0.24\textwidth}{Compatibility of late-time expansion with a $\tanh$-type template} & \parbox[t]{0.25\textwidth}{Diagnostic only; CRM was built with a $\tanh$ profile and is not an independent theorem test} & \parbox[t]{0.28\textwidth}{Model-agnostic template fit without CRM-specific priors} \\[0.6em]
\parbox[t]{0.14\textwidth}{Composition law} & \parbox[t]{0.24\textwidth}{Direct content of the projective normal form} & \parbox[t]{0.25\textwidth}{Theorem-facing empirical target once D' is assumed} & \parbox[t]{0.28\textwidth}{Compare saturation responses across redshift bins against $\mathcal{C}(x,y)=(x+y)/(1+xy)$} \\[0.6em]
\parbox[t]{0.14\textwidth}{D' defect} & \parbox[t]{0.24\textwidth}{UV selection of the projective collar} & \parbox[t]{0.25\textwidth}{Open physical derivation, not established by the QG-program table} & \parbox[t]{0.28\textwidth}{Compute or bound $\Delta_{D'}=\log R(C)-\log R(x)-\log R(y)$ in a candidate UV model} \\[0.6em]
\parbox[t]{0.14\textwidth}{Parameter values} & \parbox[t]{0.24\textwidth}{$k$, $\Phi_0$, $\gamma$ and transition sharpness} & \parbox[t]{0.25\textwidth}{CRM-/UV-specific, not determined by the theorem} & \parbox[t]{0.28\textwidth}{Derive parameters from a microscopic model and compare to CRM fits} \\[0.6em]
\parbox[t]{0.14\textwidth}{Oscillatory residuals} & \parbox[t]{0.24\textwidth}{Possible UV-sensitive late-time residual pattern} & \parbox[t]{0.25\textwidth}{Candidate signature only; not a proof of quantum origin by itself} & \parbox[t]{0.28\textwidth}{Pre-register frequency/window/null controls before interpreting BAO or SN residuals} \\
\end{tabular}
\end{ruledtabular}
\end{table*}

\textbf{Projective-collar predictions} (after D' is imposed):
\begin{enumerate}
\item If the projective boundary quotient is physically realized, the cosmological expansion history should be compatible with $\tanh$-type saturation. We emphasize that the CRM's $\Delta\chi^2 = -5.5$ result \cite{Geiger2026b} does \textit{not} constitute an independent test of the theorem, since the CRM was constructed with a $\tanh$ profile before the axioms were formulated. The theorem's empirical content is the \textit{composition law} $\mathcal{C}(x,y) = (x+y)/(1+xy)$, which is a prediction beyond any individual profile fit. An independent test would require either (a)~fitting a model-agnostic $\tanh$-template to SN+CMB+BAO data without CRM-specific priors, or (b)~verifying the composition law directly by comparing saturation responses at different redshift bins and checking consistency with the velocity-addition rule.
\item Within the CRM fit, the running curvature coupling $\beta_{\mathrm{eff}}(a)$ is reported as a monotone transition function ($\beta_{\mathrm{early}} \approx 2.8 \to \beta_{\mathrm{late}} \approx 2.0$ \cite{Geiger2026b}); its theorem-facing status is as a consistency check for Axiom~C, not as an independent derivation.
\item No alternative relaxation profile (logistic, power-law, etc.) can claim support from this theorem unless it also realizes a scale-coherent composition law and explains the projective quotient.
\end{enumerate}

\textbf{UV-dependent predictions} (discriminating between completions):
\begin{enumerate}
\item The specific values of $k$ (saturation rate), $\Phi_0$ (saturation amplitude), and $\gamma$ (coupling constant) depend on the UV theory.
\item The transition sharpness ($n$ in the sigmoidal transition) would encode a universality class if a UV-to-cosmology map is supplied.
\item Sub-Planckian corrections to the $\tanh$ profile should be treated as UV-theory-specific candidates until matched controls are defined.
\end{enumerate}

\subsection{The missing piece: why \textit{these} parameters?}

The Saturation Theorem determines the Abel/collar structure and, under D', the projective \textit{form}, but not the \textit{coefficients}. The remaining open question is:

\begin{quote}
\textit{Why does the universe have the specific values $k \approx 9.8$, $\Phi_0 \approx 0.43$, $\gamma \sim \mathcal{O}(1)\,H_0^{-2}$?}
\end{quote}

This question is analogous to the fine-structure constant problem: the \textit{existence} of electromagnetic coupling is forced by gauge symmetry, but the \textit{value} $\alpha \approx 1/137$ is not determined by the symmetry alone.

Several UV completions make specific predictions:
\begin{itemize}
\item \textbf{LQC:} $\Phi_0$ maps to the Barbero-Immirzi parameter $\gamma_{\mathrm{BI}}$, which is constrained by black hole entropy \cite{Meissner2004}. This would predict a specific value of $\Phi_0$.
\item \textbf{Asymptotic safety:} The critical exponents at the UV fixed point determine the transition sharpness $n$, providing a specific prediction for the shape of $\beta_{\mathrm{eff}}(a)$ \cite{Reuter1998}.
\item \textbf{Holography:} The de~Sitter entropy $S_{\mathrm{dS}} = \pi/(\Lambda G)$ relates $\Phi_0$ to the ratio of UV and IR scales.
\end{itemize}

A concrete test case is the $\sqrt{\pi}$ conjecture of Paper~II \cite{Geiger2026b}: the CRM predicts a MOND-like enhancement factor $\mu_{\mathrm{eff}} = \sqrt{\pi}$ arising from Gaussian normalization of the saturation partition function. Deriving this value from a specific UV completion would directly connect quantum gravity parameters to observed galactic dynamics.

\subsection{Scalar field oscillations as a QG signature}

Paper~III \cite{Geiger2026c} shows that the P\"oschl-Teller potential supports a discrete spectrum of bound states. In the cosmological context, these correspond to oscillatory corrections to the expansion rate at late times:
\begin{equation}
H^2(a) = H_{\mathrm{smooth}}^2(a)\left[1 + \epsilon \cdot e^{-\Gamma a} \cos(\omega a + \delta)\right]\,,
\end{equation}
with $\epsilon \ll 1$. These oscillations would be a candidate signature of microscopic structure in the saturation mechanism -- analogous in spirit to discrete energy levels of a quantum system. Detection in high-precision BAO or SN data would support a UV-sensitive residual channel only if the frequency window, damping scale, and null controls were fixed before fitting; it would not by itself prove the quantum origin of CRM saturation.


% ===================================================================
% 8. DISCUSSION
% ===================================================================
\section{Discussion and Outlook}
\label{sec:discussion}

\subsection{What has been achieved}

\begin{enumerate}
\item \textbf{A conditional normal-form/no-go statement:} Any cosmological theory satisfying Axioms~A--D, B', $D^\pm$, and the projective double-boundary quotient assumption D' must produce $\tanh$-type saturation. Without D', the same structural hypotheses produce the Abel/collar structure.

\item \textbf{Universality delimited:} The theorem defines a conditional Abel/collar equivalence class for responses that satisfy its gates. The QG examples currently provide motifs and mapping obligations, not a Wilsonian universality class or a shared microscopic mechanism; the projective $\tanh$ representative additionally depends on D'.

\item \textbf{CRM calibrated:} Under the stated axioms, the CRM's saturation ODE $dX/da = k(1-X^2)$ can be interpreted as the projective normal form of capacity-limited cooperative relaxation rather than only as a phenomenological fit.

\item \textbf{Ontological interpretation:} In this framework, quantum gravity may be approached not only as metric quantization, but as a theory of geometric order -- a transition from a pre-geometric quantum system to classical spacetime.
\end{enumerate}

\subsection{Relation to the CRM program}

This paper supplies the conditional axiomatic layer of the CRM programme, but it does not complete the physical arc of the series. Following the July~2026 truth audit and Route~3 reconstruction, the logical status is:
\begin{enumerate}
\item The \textit{form} of the dynamics ($\tanh$) is a mathematical theorem once the signed interior composition hypotheses hold and the projective double-boundary quotient assumption D' is added; without D' the theorem proves the Abel/collar structure.
\item The historical direct trace-coupling route of Paper~III is a failure audit; Route~3 is an action-level candidate whose full FLRW, Hamiltonian, PPN, lensing and data gates remain open.
\item The reported fit parameters constrain an effective phenomenological proxy, not yet a validated Route~3 or $f(R)$ realization.
\item The \textit{UV completion} remains open -- and must explain why the projective boundary quotient, rather than a different smooth collar, is the physical composition variable.
\end{enumerate}

\subsection{Post-Zookeeper stability-extension roadmap}

The Post-Zookeeper programme is the author's proposed future framework for classifying admissible UV response channels; it is not used in the proof of the theorem below.

\subsection{Limitations and caveats}

\begin{enumerate}
\item \textbf{Axiom D and D' are strong.} The renormalization-group composition law is restrictive, and the projective double-boundary quotient assumption is an additional mathematical and physical choice. Theories without a well-defined RG flow (some discrete approaches, multiverse scenarios) may not satisfy D; theories with a different boundary collar may satisfy D but fail D'. This distinction is essential: D gives scale coherence, D' selects the $\tanh$ representative.

\item \textbf{The proof is algebraic, not constructive.} The theorem establishes the Abel structure of the response and, under D', its projective $\tanh$ representative, but it does not construct the UV theory. The ``reparametrization'' $u(g)$ absorbs microscopic information; determining $u(g)$ from first principles requires solving the specific QG theory. A potential objection is that any bounded monotone function can trivially be written as $\tanh(\alpha\,u(g))$ for a suitable $u(g)$, making the profile shape contentless. The non-trivial content therefore resides in the \textit{composition law} and, for the $\tanh$ selection, in the projective quotient law $R(C)=R(x)R(y)$.

\item \textbf{Cosmological restriction and theory drift.} The axioms are formulated for a selected cosmological response channel and do not establish black-hole, gravitational-wave, screening or cluster phenomenology. Earlier chameleon- and Bullet-Cluster claims belong to the superseded route and are not consequences of this theorem. In the active Route~3 reconstruction, $c_T=1$ has been confirmed only on the transverse-traceless family; scalar Hamiltonian stability, FLRW response, lensing and screening remain separate gates.

\item \textbf{The continuous group classification and single-field restriction.} The proof relies on Acz\'el's representation theorem for continuous associative operations \cite{Aczel1966} and the classification of one-dimensional formal group laws \cite{Hazewinkel1978}. While these results are well-established, the physical applicability requires the assumption that the composition is \textit{one-dimensional} -- i.e., that there is a single relevant direction. This is a genuine limitation: the CRM itself (Paper~IV \cite{Geiger2026d}) introduces a vector sector alongside the scalar, and realistic modified gravity theories (e.g., Horndeski, beyond-Horndeski) generically have multiple coupled degrees of freedom. However, we note that the theorem applies to the \textit{scalar sector in isolation}: the CRM's scalar saturation channel $\Omega_\Phi(a) = \Phi_0 \tanh(k(a - a_{\mathrm{trans}}))$ is a one-dimensional response that satisfies the axioms independently of the vector sector. The multi-field extension (classifying multi-dimensional commutative group laws, including elliptic and Lubin-Tate formal groups) is an open mathematical problem that would strengthen the theorem but is not required for the scalar-sector result.

\item \textbf{Boundary-collar freedom.} Smooth boundary behavior alone does not select $\mathrm{arctanh}$. The counter-collar $\phi_\varepsilon=\operatorname{arctanh}+\varepsilon x$ preserves a smooth positive boundary collar but changes the Abel coordinate. The uniqueness selection is therefore deliberately moved to D': once the projective quotient $R=(1+x)/(1-x)$ is required to multiply under composition, $\mathrm{arctanh}=\frac12\log R$ follows immediately and uniquely up to scale.
\end{enumerate}

\subsection{The No-Go theorem as a class-internal filter}

Within the A--D+B'+$D^\pm$ framework, failure of D' means that the $\tanh$ CRM law is not selected as the projective normal form. Such a failure does not falsify the UV model; it locates it outside the projective CRM subclass or identifies a measurable D' defect. Theories satisfying Axioms~A--D together with B' and $D^\pm$ are constrained only to an Abel/collar saturation law. Theories satisfying these structural hypotheses and D' are constrained to $\tanh$ saturation in their cosmological response channel.

\subsection{Relation to other frameworks}

\textit{Swampland conjectures.} The de~Sitter swampland conjecture \cite{Obied2018} posits $|\nabla V|/V \geq c \sim \mathcal{O}(1)$ in any quantum gravity landscape, disfavoring a strict cosmological constant. The Saturation Theorem is complementary: it does not address the existence of de~Sitter vacua, but constrains the \textit{dynamical approach} to the late-time attractor. A CRM-type saturation with $\Phi_0 < \infty$ is compatible with the swampland bound, since the effective dark energy density $\rho_\Phi(a)$ evolves dynamically rather than being a fixed $\Lambda$.

\textit{EFT of Dark Energy.} The effective field theory of dark energy \cite{Gubitosi2013} parametrizes deviations from $\Lambda$CDM at the level of the perturbed action. The Saturation Theorem operates at a different level: it constrains the \textit{background} response function, not the perturbative operator basis. A future task is to express the $\tanh$ saturation in the EFT-of-DE language, mapping the CRM parameters $(k, \Phi_0, \gamma)$ to the EFT operator coefficients $\{\alpha_i(t)\}$.

\subsection{Outlook}

Three directions are immediate:

\begin{enumerate}
\item \textbf{Parameter derivation.} The next step is to derive $k$, $\Phi_0$, and $\gamma$ from a specific UV completion. The most promising candidate is the LQG/holographic intersection, where the Barbero-Immirzi parameter and the holographic entropy bound may jointly fix the saturation parameters.

\item \textbf{Multi-field extension.} Extending the theorem to theories with multiple saturation channels (e.g., scalar + vector in the CRM \cite{Geiger2026d}) requires the classification of multi-dimensional commutative group laws (the higher-dimensional analog of the Abel-equation reduction). This is mathematically richer and may yield constraints on the vector sector parameters $K_B$ and $\mathcal{F}$.

\item \textbf{Black hole interior.} If the Saturation Theorem applies to black hole interiors (with the radial coordinate replacing the scale factor) and the projective boundary quotient is the correct interior collar, it predicts $\tanh$-type singularity resolution -- consistent with LQG bounce models \cite{Ashtekar2006} and the Planck star scenario \cite{Rovelli2014}.
\end{enumerate}


% ===================================================================
% 9. CONCLUSION
% ===================================================================
\section{Conclusion}
\label{sec:conclusion}

We have proven a corrected Saturation Theorem: finite geometric capacity, cooperative reinforcement, monotone cosmic control, renormalization-group composition, signed response, and the reversible one-dimensional response-channel hypothesis $D^\pm$ reduce one-dimensional cosmological saturation to an Abel-coordinate boundary-collar structure. The $\tanh$ profile follows when the saturation boundary is additionally fixed by the projective double-boundary quotient assumption $R=(1+x)/(1-x)$.

The theorem has three consequences:

\begin{enumerate}
\item The CRM's saturation ODE $dX/da = k(1-X^2)$ can be read as the \textit{projective normal form} of capacity-limited cooperative relaxation once the boundary quotient is fixed.

\item Motifs from LQG, asymptotic safety, string theory, causal sets, and noncommutative geometry can be compared with the Abel/collar gates, but no external source-to-response homomorphism is derived here. The specific $\tanh$ representative additionally requires the projective quotient.

\item The paper offers an interpretive framework in which quantum gravity is studied not only as metric quantization, but also as the theory of geometric order -- a transition from a pre-geometric quantum system to classical spacetime.
\end{enumerate}

The remaining open problem -- deriving the specific parameters $(k, \Phi_0, \gamma)$ and the projective boundary quotient from a UV completion -- is the analog of deriving the fine-structure constant from a theory of quantum electrodynamics. The theorem constrains the \textit{normal form} of any theory satisfying the stated hypotheses once the boundary quotient is fixed.

\begin{quote}
\textit{The saturation profile is constrained to the $\tanh$ type by scale coherence plus the projective boundary quotient; bare smooth boundary regularity gives the wider Abel-collar class.}
\end{quote}


% ===================================================================
% AI DISCLOSURE
% ===================================================================
\begin{acknowledgments}
This work was conducted as independent research without institutional funding. Public scientific-computing infrastructure was used throughout the CRM program.

\textit{Tool use.} AI-based tools provided non-authorial process support for literature organization, drafting and code assistance, language editing, consistency checks, and \LaTeX/PDF quality assurance. Source selection, mathematical arguments, physical hypotheses, axiomatic choices, scientific interpretation, and final decisions were reviewed and retained by the author. The tools are not authors or co-authors, are not acknowledged as contributors, and did not serve as sources of truth or independent validation. The author takes full responsibility for the content.

\textit{Funding information.} This research received no external funding.

\textit{Data availability.} This paper is purely theoretical. The public CRM repository with analysis scripts, data products, and paper source files is available at \url{https://github.com/research-line/crm-cosmology}; the CRM I--IV Zenodo concept DOI is \href{https://doi.org/10.5281/zenodo.18728935}{10.5281/zenodo.18728935}.
\end{acknowledgments}


% ===================================================================
% BIBLIOGRAPHY
% ===================================================================
\begin{thebibliography}{50}

\bibitem{Geiger2026}
L.~Geiger (2026).
Game-Theoretic Cosmology and the Curvature Relaxation Model: Nash Equilibria Between Null Space and Spacetime Bubble.
Companion paper (Paper~I).
\url{https://github.com/research-line/crm-cosmology/tree/main/papers}.
\href{https://doi.org/10.5281/zenodo.18728935}{doi:10.5281/zenodo.18728935}.

\bibitem{Geiger2026b}
L.~Geiger (2026).
Eliminating the Dark Sector: Unifying the Curvature Relaxation Model with MOND.
Companion paper (Paper~II).
\href{https://doi.org/10.5281/zenodo.18728935}{doi:10.5281/zenodo.18728935}.

\bibitem{Geiger2026c}
L.~Geiger (2026).
Microscopic Foundations of the Curvature Relaxation Model: From Quantum Geometry to Macroscopic Saturation.
Companion paper (Paper~III).
\href{https://doi.org/10.5281/zenodo.18728935}{doi:10.5281/zenodo.18728935}.

\bibitem{Geiger2026d}
L.~Geiger (2026).
The Galactic--Cosmological Nexus: Connecting CRM to MOND Dynamics via Curvature Saturation.
Companion paper (Paper~IV).
\href{https://doi.org/10.5281/zenodo.18728935}{doi:10.5281/zenodo.18728935}.

\bibitem{Donoghue1994}
J.~F.~Donoghue (1994).
General relativity as an effective field theory: The leading quantum corrections.
\textit{Physical Review D}, 50(6), 3874.
DOI: 10.1103/PhysRevD.50.3874.
arXiv:gr-qc/9405057.

\bibitem{Eichhorn2020}
A.~Eichhorn,
\textit{Asymptotically safe gravity},
arXiv:2003.00044 (2020).

\bibitem{EichhornHeld2018}
A.~Eichhorn and A.~Held,
\textit{Top mass from asymptotic safety},
Phys.\ Lett.\ B \textbf{777}, 217--221 (2018).
DOI: 10.1016/j.physletb.2017.12.040.
arXiv:1707.01107.

\bibitem{Weinberg1979}
S.~Weinberg (1979).
Ultraviolet divergences in quantum theories of gravitation.
In S.~W.~Hawking \& W.~Israel (Eds.), \textit{General Relativity: An Einstein Centenary Survey}, pp.\ 790--831.
Cambridge University Press.

\bibitem{Reuter1998}
M.~Reuter (1998).
Nonperturbative evolution equation for quantum gravity.
\textit{Physical Review D}, 57(2), 971.
DOI: 10.1103/PhysRevD.57.971.
arXiv:hep-th/9605030.

\bibitem{Percacci2017}
R.~Percacci (2017).
\textit{An Introduction to Covariant Quantum Gravity and Asymptotic Safety}.
World Scientific.

\bibitem{Polchinski1998}
J.~Polchinski (1998).
\textit{String Theory}, Vols.~1--2.
Cambridge University Press.

\bibitem{Rovelli2004}
C.~Rovelli (2004).
\textit{Quantum Gravity}.
Cambridge University Press.

\bibitem{Thiemann2007}
T.~Thiemann (2007).
\textit{Modern Canonical Quantum General Relativity}.
Cambridge University Press.

\bibitem{Maldacena1999}
J.~Maldacena (1998).
The large-$N$ limit of superconformal field theories and supergravity.
\textit{Adv.\ Theor.\ Math.\ Phys.}\ \textbf{2}, 231--252.
[Reprinted in \textit{Int.\ J.\ Theor.\ Phys.}\ \textbf{38}, 1113--1133 (1999).]
DOI: 10.4310/ATMP.1998.v2.n2.a1.
arXiv:hep-th/9711200.

\bibitem{Sorkin2003}
R.~D.~Sorkin (2005).
Causal sets: Discrete gravity.
In \textit{Lectures on Quantum Gravity}, pp.\ 305--327.
Springer.
DOI: 10.1007/0-387-24992-3\_7.
arXiv:gr-qc/0309009.

\bibitem{Surya2019}
S.~Surya (2019).
The causal set approach to quantum gravity.
\textit{Living Reviews in Relativity}, 22, 5.
DOI: 10.1007/s41114-019-0023-1.

\bibitem{Connes1994}
A.~Connes (1994).
\textit{Noncommutative Geometry}.
Academic Press.

\bibitem{Bekenstein1981}
J.~D.~Bekenstein (1981).
Universal upper bound on the entropy-to-energy ratio for bounded systems.
\textit{Physical Review D}, 23(2), 287--298.
DOI: 10.1103/PhysRevD.23.287.

\bibitem{Almheiri2015}
A.~Almheiri, X.~Dong, \& D.~Harlow (2015).
Bulk locality and quantum error correction in AdS/CFT.
\textit{Journal of High Energy Physics}, 2015, 163.
DOI: 10.1007/JHEP04(2015)163.

\bibitem{Aczel1966}
J.~Acz\'el (1966).
\textit{Lectures on Functional Equations and Their Applications}.
Academic Press, New York.

\bibitem{MesiarovaZemankova2015}
A.~Mesiarov\'a-Zem\'ankov\'a (2015).
Characterization of uninorms with continuous underlying t-norm and t-conorm by means of an extended ordinal sum.
arXiv:1506.07820.
DOI: 10.48550/arXiv.1506.07820.

\bibitem{Grabisch2009}
M.~Grabisch, J.-L.~Marichal, R.~Mesiar, \& E.~Pap (2009).
Aggregation on bipolar scales.
In \textit{Proceedings of the 30th Linz Seminar on Fuzzy Set Theory}, 46--51.

\bibitem{Hazewinkel1978}
M.~Hazewinkel (1978).
\textit{Formal Groups and Applications}.
Academic Press.

\bibitem{Ashtekar2006}
A.~Ashtekar, T.~Pawlowski, \& P.~Singh (2006).
Quantum nature of the Big Bang.
\textit{Physical Review Letters}, 96(14), 141301.
DOI: 10.1103/PhysRevLett.96.141301.
arXiv:gr-qc/0602086.

\bibitem{Bojowald2008}
M.~Bojowald (2008).
Loop quantum cosmology.
\textit{Living Reviews in Relativity}, 11, 4.
DOI: 10.12942/lrr-2008-4.

\bibitem{Bonanno2002}
A.~Bonanno \& M.~Reuter (2002).
Cosmology of the Planck era from a renormalization group for quantum gravity.
\textit{Physical Review D}, 65(4), 043508.
DOI: 10.1103/PhysRevD.65.043508.
arXiv:hep-th/0106133.

\bibitem{Ryu2006}
S.~Ryu \& T.~Takayanagi (2006).
Holographic derivation of entanglement entropy from the anti-de~Sitter space/conformal field theory correspondence.
\textit{Physical Review Letters}, 96(18), 181602.
DOI: 10.1103/PhysRevLett.96.181602.
arXiv:hep-th/0603001.

\bibitem{deBoer2000}
J.~de Boer, E.~P.~Verlinde, \& H.~L.~Verlinde (2000).
On the holographic renormalization group.
\textit{Journal of High Energy Physics}, 2000(08), 003.
DOI: 10.1088/1126-6708/2000/08/003.
arXiv:hep-th/9912012.


\bibitem{Yang1952}
C.~N.~Yang (1952).
The spontaneous magnetization of a two-dimensional Ising model.
\textit{Physical Review}, 85, 808--816.
DOI: 10.1103/PhysRev.85.808.

\bibitem{WilsonKogut1974}
K.~G.~Wilson \& J.~Kogut (1974).
The renormalization group and the $\epsilon$ expansion.
\textit{Physics Reports}, 12(2), 75--199.
DOI: 10.1016/0370-1573(74)90023-4.

\bibitem{Meissner2004}
K.~A.~Meissner (2004).
Black-hole entropy in Loop Quantum Gravity.
\textit{Classical and Quantum Gravity}, 21(22), 5245.
DOI: 10.1088/0264-9381/21/22/015.
arXiv:gr-qc/0407052.

\bibitem{Rovelli2014}
C.~Rovelli \& F.~Vidotto (2014).
Planck stars.
\textit{International Journal of Modern Physics D}, 23(12), 1442026.
DOI: 10.1142/S0218271814420267.
arXiv:1401.6562.

\bibitem{PDG2024}
S.~Navas et~al. (Particle Data Group) (2024).
Review of Particle Physics.
\textit{Physical Review D}, 110, 030001.
DOI: 10.1103/PhysRevD.110.030001.

\bibitem{Obied2018}
G.~Obied, H.~Ooguri, L.~Spodyneiko, \& C.~Vafa (2018).
De Sitter Space and the Swampland.
arXiv:1806.08362.

\bibitem{Gubitosi2013}
G.~Gubitosi, F.~Piazza, \& F.~Vernizzi (2013).
The Effective Field Theory of Dark Energy.
\textit{Journal of Cosmology and Astroparticle Physics}, 2013(02), 032.
DOI: 10.1088/1475-7516/2013/02/032.
arXiv:1210.0201.

\bibitem{Ambjorn2005}
J.~Ambj\o{}rn, J.~Jurkiewicz, \& R.~Loll (2005).
The spectral dimension of the universe is scale dependent.
\textit{Physical Review Letters}, 95(17), 171301.
DOI: 10.1103/PhysRevLett.95.171301.
arXiv:hep-th/0505113.

\bibitem{Loll2020}
R.~Loll (2020).
Quantum gravity from causal dynamical triangulations: a review.
\textit{Classical and Quantum Gravity}, 37(1), 013002.
DOI: 10.1088/1361-6382/ab57c7.
arXiv:1905.08669.

\bibitem{HartnollYang2025}
S.~A.~Hartnoll \& M.~Yang (2025).
The conformal primon gas at the end of time.
\textit{Journal of High Energy Physics}, 2025(07), 281.
DOI: 10.1007/JHEP07(2025)281.
arXiv:2502.02661.

\bibitem{Julia1990}
B.~L.~Julia (1990).
Statistical theory of numbers.
In J.~M.~Luck, P.~Moussa, \& M.~Waldschmidt (Eds.), \textit{Number Theory and Physics},
Springer Proceedings in Physics, Vol.~47, pp.\ 276--293.
Springer-Verlag, Berlin.

\bibitem{PapadopoulosTroyanov2008}
A.~Papadopoulos \& M.~Troyanov (2008).
Harmonic symmetrization of convex sets and of Finsler structures, with applications to Hilbert geometry.
arXiv:0807.0335.

\bibitem{RainioVuorinen2023}
O.~Rainio \& M.~Vuorinen (2023).
Hilbert metric in the unit ball.
arXiv:2303.03753.

\bibitem{FendleyLudwigSaleur1995}
P.~Fendley, A.~W.~W.~Ludwig, \& H.~Saleur (1995).
Exact conductance through point contacts in the $\nu=1/3$ fractional quantum Hall effect.
\textit{Physical Review Letters}, 74, 3005--3008.
DOI: 10.1103/PhysRevLett.74.3005. arXiv:cond-mat/9408068.

\bibitem{GhoshalZamolodchikov1994}
S.~Ghoshal \& A.~B.~Zamolodchikov (1994).
Boundary S matrix and boundary state in two-dimensional integrable quantum field theory.
\textit{International Journal of Modern Physics A}, 9, 3841--3886.
DOI: 10.1142/S0217751X94001552. arXiv:hep-th/9306002.

\bibitem{KostrykinSchrader2001}
V.~Kostrykin \& R.~Schrader (2001).
The generalized star product and the factorization of scattering matrices on graphs.
\textit{Journal of Mathematical Physics}, 42, 1563--1598.
DOI: 10.1063/1.1354641. arXiv:math-ph/0008022.

\bibitem{FendleySaleurZamolodchikov1993}
P.~Fendley, H.~Saleur, \& A.~B.~Zamolodchikov (1993).
Massless flows II: The exact S-matrix approach.
\textit{International Journal of Modern Physics A}, 8, 5751--5778.
DOI: 10.1142/S0217751X93002277. arXiv:hep-th/9304051.

\bibitem{HeJiangLiu2026}
Y.~He, Y.~Jiang, \& Y.~Liu (2026).
Fusion of integrable defects and the defect $g$-function.
arXiv preprint arXiv:2605.20688v1.

\bibitem{DelfinoMussardoSimonetti1994b}
G.~Delfino, G.~Mussardo, \& P.~Simonetti (1994).
Scattering theory and correlation functions in statistical models with a line of defect.
\textit{Nuclear Physics B}, 432, 518--550.
DOI: 10.1016/0550-3213(94)90032-9. arXiv:hep-th/9409076.

\bibitem{Selberg1956}
A.~Selberg (1956).
Harmonic analysis and discontinuous groups in weakly symmetric Riemannian spaces with applications to Dirichlet series.
\textit{Journal of the Indian Mathematical Society}, 20, 47--87.
DOI: 10.18311/JIMS/1956/16985.

\bibitem{Ungar2007}
A.~A.~Ungar (2007).
Einstein's velocity addition law and its hyperbolic geometry.
\textit{Computers \& Mathematics with Applications}, 53(8), 1228--1250.
DOI: 10.1016/j.camwa.2006.05.028.

\end{thebibliography}

\end{document}
